\documentclass{article}
\usepackage{amssymb,amsmath,amsthm}

% Referencing items in an enumerate environment
\usepackage{enumitem}

% Flag toggling functionality
\usepackage{etoolbox}

% For automatic spacing tweaks to avoid overfull boxes
\usepackage{microtype}

\usepackage[style=numeric,sorting=none]{biblatex}
\addbibresource{\jobname.bib}

% Conditionals
%-----------------------------------------------------------------------------%

% Draft status
\newtoggle{draft}

% Title page
\newtoggle{titlepage}

% Contents page
\newtoggle{contentspage}

% Colour inversion
\newtoggle{dark}
\newtoggle{colour}

% Toggle
\togglefalse{draft}
\togglefalse{colour}
\togglefalse{dark}
\toggletrue{titlepage}
\toggletrue{contentspage}

%-----------------------------------------------------------------------------%

% Title page needs graphicx to use \includegraphics
\iftoggle{titlepage}{%
	\usepackage{graphicx}
}{}

% Watermark
\iftoggle{draft}{%
	\usepackage{draftwatermark}

	\SetWatermarkScale{5}

	% Set font to supported size. If different size required, use a suitable
	% package to avoid warnings.
	\SetWatermarkFontSize{24.88pt}
}

% Document colouring
\iftoggle{colour}{%
	\usepackage{xcolor}

	\iftoggle{dark}{%
		\pagecolor[RGB]{40,40,40}
		\color[RGB]{235,219,178}
	}{%
		\pagecolor[RGB]{251,241,199}
		\color[RGB]{60,56,54}
	}
}{}

% Nomenclature
\newcommand\name{z0}
\newcommand\paper{paper}

\newcommand\newProtocol[2]{%
	\expandafter\newcommand\csname f#1\endcsname[1]{%
		f_\textup{\texttt{#2}}^{##1}
	}%
	\expandafter\newcommand\csname pi#1\endcsname[1]{%
		\pi_\textup{\texttt{#2}}^{##1}
	}%
}

% Mathematical notation
\newcommand{\share}[2]{\left[#1\right]_{#2}}
\newcommand{\eqd}[0]{\stackrel{d}{=}}
\newcommand{\powerset}{\raisebox{.15\baselineskip}{\Large\ensuremath{\wp}}}
\newcommand{\set}[2]{\left\{ #1 \middle| #2 \right\}}
\newcommand{\seq}[1]{\left[#1\right]}
\newcommand{\seqZ}[1]{\left[#1\right]_0}

% Mathematical environments
\newtheorem{theorem}{Theorem}
\newtheorem{corollary}{Corollary}
\newtheorem{lemma}{Lemma}
\newtheorem{claim}{Claim}
\newtheorem{definition}{Definition}
\newtheorem{assumption}{Assumption}

\theoremstyle{remark}
\newtheorem{remark}{Remark}

% Protocols
\newProtocol{Open}{OPEN}
\newProtocol{SOpen}{SO}
\newProtocol{DOpen}{DO}
\newProtocol{SDOpen}{SDO}
\newProtocol{BRNG}{BRNG}
\newProtocol{RNG}{RNG}
\newProtocol{RZG}{RZG}
\newProtocol{RKPG}{RKPG}
\newProtocol{MulOpen}{MO}
\newProtocol{Mul}{MUL}
\newProtocol{Inv}{INV}
\newProtocol{Sign}{SIGN}

% Common notation
\newcommand{\N}{\mathbb{N}}
\newcommand{\F}{\mathbb{F}}
\newcommand{\G}{\mathbb{G}}
\renewcommand{\P}{\mathbb{P}}
\newcommand{\ind}{\mathcal{I}}
\newcommand{\ic}{\mathcal{I}_c}
\newcommand{\ih}{\mathcal{I}_h}

\begin{document}

\iftoggle{titlepage}{%
	\begin{titlepage}
		\begin{center}
			\vspace*{1cm}

			\Huge \textbf{RenVM Secure Multiparty Computation}

			\vspace{0.5cm}

			%\LARGE Subtitle

			\vspace{1.5cm}

			\large\textbf{Ross Pure, Zian-Loong Wang}

			\vspace{4.0cm}

			\includegraphics[width=0.4\textwidth]{ren_logo_black.png}
		\end{center}
	\end{titlepage}
}

\iftoggle{contentspage}{%
	\tableofcontents
	\newpage
}

\section{Introduction}

A threshold ECDSA signature scheme allow groups of $n$ players to jointly
produce an ECDSA signature without exposing the corresponding ECDSA private
key, and no subset of less than $t+1$ players is able to produce a signature.
Generally, such schemes also allow $n$ players to jointly generate an ECDSA
public/private key pair, without the need for a dealer, and where the public
key is known by everyone.

Interest in threshold ECDSA schemes has increased recently, with the
popularisation of blockchains (such as Bitcoin and Ethereum) that make use of
ECDSA signatures to custody and spend cryptocurrencies. Using threshold ECDSA
schemes, groups of $n$ players can take joint custody of cryptocurrencies, and
only spend these funds when at least $t+1$ players agree that this should be
done.

The most recent threshold DSA scheme, which is general enough to support ECDSA
specifically, was introduced by Gennaro et al.~\cite{gg18}. This scheme allows
$n \ge t+1$ players to generate public/private keys and produce signatures.
This work is primarily an extension of~\cite{ggn16}, improving its efficiency
and reducing its complexity. Threshold DSA schemes were previously introduced
by~\cite{gjkr96}, and~\cite{mr04}. The main advantage of~\cite{ggn16}
and~\cite{gg18} over these schemes is threshold-optimality (the scheme is
secure against $t$ adversaries, and only requires participation from $t+1$
players).

In the context of large-scale globally distributed Byzantine networks, where
participants are distributed all around the world, can be anyone, are
constantly changing, and latency is large and highly variant, robustness
against offline participants is critical to the practical usability of any
protocol. In~\cite{gg18}, the use of additive secret sharing, inspired by
SPDZ~\cite{dkl+12}, introduces moments in the protocol where failure to
participate by one player causes the protocol to halt (and it needs to be
repeated from a prior round).  This might be usable for small groups, but it is
impractical for large globally distributed groups where honest players are
likely to experience unexpected downtime or network outages. In~\cite{ggn16},
the need for trusted setup of an additively homomorphic encryption scheme (the
removal of which was the core improvement made by~\cite{gg18}) prevents it from
being efficient with a large number of players in a high-latency network.

In~\cite{gjkr96}, a scheme is presented that is robust against malicious
adversaries, but makes extensive use of secure broadcasts and rounds of
complaining against misbehaving players. This introduces extra rounds of
communication, and more constraints on network timing, which make it
inappropriate for large-scale globally distributed networks.  On the other
hand~\cite{mr04} requires $O(t)$ rounds of communication using Paillier’s
encryption with a modulus $N = O(q^{3t-1})$, making it unsuitable for a large
number of participants.

In this paper, we present a threshold ECDSA scheme, based on~\cite{gjkr96},
that is practical for use in large-scale globally distributed networks. It is
robustly secure against $t$ malicious adversaries, assuming $n = 3t + 1$
players, and requires fewer rounds than~\cite{gjkr96}. During both ECDSA key
generation and signing, up to $t$ players can go offline at the beginning,
middle, or end of a round, and the protocols will complete successfully without
the need to go back repeat from a prior round.

\subsection{Contributions}

Our threshold ECDSA scheme is based on~\cite{gjkr96} but replaces the protocols
that generate the random sharings (random number generation, random zero
generation, and random key-pair generation) with ones that don't require secure
broadcasts and can be proven in Canetti's security framework~\cite{c00}.
Briefly, these new protocols are as follows.
\begin{itemize}
	\item (Biased RNG $\piBRNG{n,k}$): This is similar to the RNG protocol
		in~\cite{p91a}, but we call it biased because of the ability of a
		rushing adversary to influence the sharing (but not the underlying
		random number). The exact protocol is similar to the one
		in~\cite{gjkr96} but secure broadcasting is replaced with a Byzantine
		fault tolerant consensus protocol.
	\item (RNG $\piRNG{n,k}$): This is the unbiased version of $\piBRNG{n,k}$.
		It is built using $\piBRNG{n,k}$. Because the shares of honest players
		cannot be biased, we are able to prove its security in Canetti’s
		framework.
	\item (Random zero generation $\piRZG{n,k}$): This is the same as
		$\piRNG{n,k}$ except that instead of generating shares of a random
		value, it generates shares of the additive identity (zero element) of
		the field. In~\cite{gjkr96} this is achieved using a variation of their
		robust RNG protocol. However, in this \paper{} we will use a protocol
		very similar to our $\piRNG{n,k}$ to again avoid the biasing problem.
	\item (Random key-pair generation $\piRKPG{n,k}$): This is a distributed
		key-pair generation protocol built using $\piRNG{n,k}$. A distributed
		key-pair generation protocol is also presented in~\cite{gjkr07}, but it
		requires extra rounds of secure broadcasting to remove biasing of the
		public key.
\end{itemize}

Additionally, we also present proofs of security for open protocols; these
protocols are the familiar opening of a shared secret. The proofs are given for
Canetti's security framework, as we want to prove that the ECDSA signature
algorithm is secure in this framework, and opens are needed for multiplying
shared secrets and also inverting a shared secret. An added benefit of having
security proofs of our subprotocols in Canetti's framework is that from the
composability property, they can be readily used for other applications as
well. However, we will see that for producing ECDSA signatures, we get the
(weaker) security property proven in~\cite{gjkr96} for free as the protocols
are almost the same.

\subsection{Organisation}

The rest of the \paper{} is structured as follows. First, background material
including probability theory, Shamir secret sharing and computational
assumptions will be presented. Next, the security framework to be used is
described. The ECDSA signature algorithm will then be introduced, and how it
will be achieved in SMPC, making comparisons to previous work. The rest of the
paper is dedicated to the SMPC protocols and their security. Initially, the
more primitive protocols such as opening shared secrets and random number
generation are described. Next, higher level protocols like multiplication and
field inversion are constructed using these primitives. Finally, we will
present the formal theorems for the security of our protocol for computing
ECDSA signatures.

\section{Background}

In this section we outline some background material that will be used
throughout the \paper{}. This includes basic probability theory definitions and
theorems, Shamir Secret Sharing, cryptographic assumptions and also some
notation.

\subsection{Notation}

Here we outline some common notation that will be used in this \paper{}. The
natural numbers are denoted by the set $\N$, where $0 \notin \N$; i.e.\ $\N =
\{1, 2, \ldots\}$. We will use $\F$ to denote an arbitrary field (in this
\paper{} we will only be concerned with finite fields). When talking about
field elements, it is understood that $0$ and $1$ represent respectively the
additive and multiplicative identity of the field. An anonymous probability
measure is represented by $\P$; this will appear in cases where a probability
measure is being used where one has not been explicitly defined, and the
definition of it should be clear from the context. The number of elements in a
finite set $A$ will be denoted $|A|$. Generally, we will denote index sets as
$\mathcal{I}$ or its non calligraphic form $I$. We also state the following
definitions. For a given index set $\mathcal{I}$, we define the notation
${(x_i)}_{i \in \mathcal{I}}$ to be the $|\mathcal{I}|$-tuple of elements
$(x_{i_1}, x_{i_2}, \ldots, x_{i_n})$, where $\{i_1, \ldots, i_n\} =
\mathcal{I}$. In all cases in this \paper{}, there will exist some order on
$\mathcal{I}$ (almost always we will have $\mathcal{I} \subset \F$ for some
finite field $\F$), and so in the above we will assume that $i_1 < i_2 < \cdots
< i_n$. The empty set will be denoted $\emptyset$.

\begin{definition}
	Let a set $A$ be given. Denote the \textit{power set} of $A$, i.e.\ the set
	of all subsets of $A$, by $\powerset(A)$.
\end{definition}

\begin{definition}
	Let the set of consecutive natural numbers from 1 to $n$ inclusive, i.e.\
	$1, 2, \ldots, n$, be denoted by $\seq{n}$.
\end{definition}

\begin{definition}
	Let the set of consecutive natural numbers from 0 to $n$ inclusive, i.e.\
	$0, 1, \ldots, n$, be denoted by $\seqZ{n}$.
\end{definition}

\subsection{General Definitions}

A \textit{negligible} function is one that approaches zero very fast. This is
made precise in the following definition.
\begin{definition}[Negligible function]
	A function $f: \N \to \mathbb{R}$ is negligible if for all $c \in \N$ there
	exists some $N \in \N$ such that for all $x \ge N$
	\begin{align*}
		|f(x)| < \frac{1}{x^c}.
	\end{align*}
	In this case we say that $f$ \textbf{is a negligible function}, or that $f$
	\textbf{is negligible}.
\end{definition}

\subsection{Probability}

Many of the security definitions and consequently also the proofs of security
are formulated in terms of probability theory. We therefore collect some key
definitions and theorems that we will use in this \paper{} here. For more
detailed on the underlying measure and probability theory, some key definitions
and theorems are collected in Appendix~\ref{app:prob} for convenience.

We first define what it means for two random variables to have equal
distributions. The idea is that the possible outcomes for the random variables
should have the same probabilities.

\begin{definition}
	Let $(\Omega_1, \Sigma_1, \P_1)$ and $(\Omega_2, \Sigma_2, \P_2)$ be
	probability spaces, $(E, \mathcal{E})$ a measurable space, and $X_1:
	\Omega_1 \to E$ and $X_2: \Omega_2 \to E$ be two random variables. We say
	that $X_1$ and $X_2$ are \textbf{identically distributed} if $\forall A \in
	\mathcal{E}$ we have
	\begin{align*}
		\P_1\left(\set{\omega \in \Omega_1}{X_1(\omega) \in A}\right) =
		\P_2\left(\set{\omega \in \Omega_2}{X_2(\omega) \in A}\right).
	\end{align*}
	In this case we write $X_1 \eqd X_2$.
\end{definition}

Here we present a result that will be useful in the proofs of security. This
formalises the notion that adding a uniformly random element to some other
element results in an element that is itself uniformly random. This is
presented abstractly in the following theorem, and then the result of interest
follows as a corollary. The proof can be found in Appendix~\ref{app:unif}.

\begin{theorem}\label{thm:unif}
	Let $(\Omega, \Sigma, \P)$ be a probability space, $E$ a finite set with
	$n$ elements, and $f_i: \Omega \to E$ for $1 \le i \le n$ be functions that
	are measurable in the context of the measurable spaces $(\Omega, \Sigma)$
	and $(E, \powerset(E))$ and satisfy the property
	\begin{align*}
		f_i(\omega) \ne f_j(\omega)
		\; \forall \omega \in \Omega
		\; \forall i \ne j\;\text{where}\;i, j \in \{1, \ldots, n\}.
	\end{align*}
	Let $F = \{f_1, \ldots, f_n\}$ and consider the probability space $(F,
	\powerset(F), U_F)$. Construct the product probability space $(\Omega
	\times F, \Sigma \times \powerset(F), \mu)$ where $\mu$ is the appropriate
	product measure. Define the random variable
	\begin{align*}
		X: \Omega \times F &\to E\\
		(\omega, f) &\mapsto f(\omega)
	\end{align*}
	Then $X$ is uniformly distributed on $E$.
\end{theorem}

\begin{corollary}\label{cor:unif}
	Let $(\F^n, \powerset(\F^n), \P)$ be a probability space, where $\F$ is a
	finite field. Let $x, y \in \F^n$ be random such that $x$ is distributed
	according to $\P$, and $y$ is uniformly distributed and independent from
	$x$. Then the random variable $x + y$ is uniformly distributed, where when
	$n > 1$ the addition is element-wise.
\end{corollary}

\begin{proof}
	This follows from Theorem~\ref{thm:unif} by setting $(\Omega, \Sigma, \P) =
	(\F^n, \powerset(\F^n), \P)$, $E = \F^n$, and defining $f_i$ by
	$f_i(\omega) = \omega + i$ for all $i \in \F^n$.
\end{proof}

Here we formalise another result that will be used in the proofs of security,
to ease their exposition. The result is simple: if two random variables have
equal distributions, then applying some deterministic function $f$ to these
random variables results in two new random variables that also have equal
distributions.

\begin{theorem}\label{thm:eqd_det}
	Let $(\Omega_1, \Sigma_1, \P_1)$ and $(\Omega_2, \Sigma_2, \P_2)$ be
	probability spaces, $(E, \mathcal{E})$ be a measurable space, and $X_1:
	\Omega_1 \to E$ and $X_2: \Omega_2 \to E$ be random variables. Let $(F,
	\mathcal{F})$ be a measurable space and $f: E \to F$ be a measurable
	function. Then if $X_1 \eqd X_2$, it follows that $f(X_1) \eqd f(X_2)$.
\end{theorem}

\begin{proof}
	Define the random variables $Y_1 = f \circ X_1$ and $Y_2 = f \circ X_2$ and
	let $A \in \mathcal{F}$ be arbitrary. We need to show that
	$\P_1(Y_1^{-1}(A)) = \P_2(Y_2^{-1}(A))$. But by assumption
	$\P_1(X_1^{-1}(B)) = \P_2(X_2^{-1}(B))$ for any $B \in \mathcal{E}$, and in
	particular for $B = f^{-1}(A)$.
\end{proof}

\subsection{Secret Sharing}\label{sec:sss}

\newcommand{\spl}[1][k]{\vartheta_{#1}}

We will consider the secret sharing technique of Shamir~\cite{shamir_1979}.
Shamir secret sharing is a technique that allows for the distribution of a
secret to $n$ players such that any subset of $k$ or more players can
reconstruct the secret, but any less than $k$ can learn nothing about the
secret. The parameters $n, k \in \N$ (where $n \ge k$) can be chosen as
required.

Any secret in a field can be shared. Let $(\F, +, \;\cdot\;)$ be a finite field
and $n, k \in \N$ be given where $n \ge k$. Let the $n$ players be
${\{P_i\}}_{i \in \mathcal{I}}$ where $\mathcal{I} \subset \F \setminus \{0\}$
and $|\mathcal{I}| = n$. We can ``share'' a secret $s \in \F$ by first choosing
uniformly random $c_1, \ldots, c_{k-1} \in \F$ and constructing the polynomial
\begin{align*}
	p(x) = s + \sum_{j=1}^{k-1} c_j x^j.
\end{align*}
Player $P_i$ is given a ``share'' $s_i = p(i)$. Note that $p(0) = s$ which is
why we should ensure that $0 \notin \mathcal{I}$.

Any subset of players ${\{P_i\}}_{i \in \mathcal{R}}$ such that $\mathcal{R}
\subset \mathcal{I}$ and $|\mathcal{R}| \ge k$ can reconstruct the secret $s$
by interpolating to reconstruct the original polynomial as follows:
\begin{align*}
	p(x)
	=
	\sum_{i \in \mathcal{R}}
		s_i
		\prod_{\substack{j \in \mathcal{R}\\j \ne i}}
			\frac{x - x_j}{x_i - x_j}.
\end{align*}
The secret is simply $p(0)$ and so the players can find $s$ using the above;
\begin{align*}
	s = \sum_{i \in \mathcal{R}}
		s_i
		\prod_{\substack{j \in \mathcal{R}\\j \ne i}}
			\frac{x_j}{x_j - x_i}.
\end{align*}
Define the function that takes the random coefficients to resulting shares by
\begin{align*}
	\spl:
	\powerset(\F) \times \F^k &\to \F^{|\mathcal{I}|}\\
	\mathcal{I}, (c_0, \ldots, c_{k-1})
	&\mapsto
	{\left(\sum_{i=0}^{k-1} j^i c_i\right)}_{j \in \mathcal{I}},
\end{align*}
where naturally the secret is $c_0$.

We now present a result that tells us how a subset of shares is distributed
given that we know how the secret is distributed. This result will be useful
later on in our security proofs. A similar result can be found in~\cite{al17}.
The proof of this result can be found in Appendix~\ref{app:shamir}.

\begin{theorem}\label{thm:shareDist}
	Let $x_1, \ldots, x_n$ be a $k$-sharing of some $x \in \F$, where $x$ has
	distribution determined by the probability space $(\F, \powerset(\F),
	\mu)$. Suppose that the sharing is a uniformly random one; that the
	remaining coefficients $c_1, \ldots, c_{k-1}$ that define the sharing are
	all uniformly and independently distributed. Let $I \subset \mathcal{I}$
	with $|I| = t$. Precisely, the set ${(x_i)}_{i \in I}$ is the random
	variable
	\begin{align}\label{eq:shareDist}
		\begin{aligned}
			X: \F^k &\to \F^t\\
			(x, c_1, \ldots, c_{k-1}) &\mapsto
			{\left(x + \sum_{j=1}^{k-1} c_j i^j\right)}_{i \in I}
		\end{aligned}
	\end{align}
	which has the associated product probability measure that we will denote by
	$\P_t$. Then this measure satisfies
	\begin{align*}
		\P_t(X = z) =
		\begin{cases}
			\mu(\{y\})|\F|^{-k+1} & t \ge k\\
			|\F|^{-t} & t < k
		\end{cases}
	\end{align*}
	for all $z \in \F^t$ that form part of some consistent $k$-sharing, where
	in the former case $y \in \F$ is determined uniquely from $z$.
\end{theorem}

Often, we will need to consider what the distribution of the shares are, given
some subset that has been fixed. The following definition introduces a notation
for the set of possible sharings given the fixed subset.

\newcommand{\shareRestrict}[3]{S_{#1, #2, #3}}

\begin{definition}
	Let $t, k, n \in \N$ with $t < k \le n$. Let $\mathcal{I} \in \F^n$ be an
	index set with $|\mathcal{I}| = n$ and let $I \subset \mathcal{I}$ be such
	that $|I| = t$. Let the shares $x = {(x_i)}_{i \in I} \in \F^t$ be fixed.
	Then we define the set $\shareRestrict{x}{\mathcal{I}}{I}$ to be the set of
	possible collections of $n - t$ shares that extend $x$ to a consistent
	$k$-sharing that has $n$ shares of some field element. Precisely,
	$\shareRestrict{x}{\mathcal{I}}{I}$ is the set of all values ${(y_i)}_{i
	\in \mathcal{I} \setminus I}$ in $\F^{n-t}$ such that there exists $(c_0,
	\ldots, c_{k-1}) \in \F^k$ for which the following hold:
	\begin{align*}
		x_i &= \sum_{j=0}^{k-1} c_j i^j \quad
		\forall i \in I,\\
		y_i &= \sum_{j=0}^{k-1} c_j i^j \quad
		\forall i \in \mathcal{I} \setminus I.
	\end{align*}
\end{definition}

An easy property of $\shareRestrict{x}{\mathcal{I}}{I}$ is that it has
$|\F|^{k-t}$ elements. This is summarised in the following theorem.

\begin{theorem}\label{thm:restrictSize}
	Let $t, k, n \in \N$ with $t < k \le n$. Let $\mathcal{I} \in \F^n$ be an
	index set with $|\mathcal{I}| = n$ and let $I \subset \mathcal{I}$ be such
	that $|I| = t$. Let the shares $x = {(x_i)}_{i \in I} \in \F^t$ be fixed.
	Then
	\begin{align*}
		|\shareRestrict{x}{\mathcal{I}}{I}| = |\F|^{k-t}.
	\end{align*}
\end{theorem}

\begin{proof}
	We know that if we have $k$ fixed shares $y_1, \ldots, y_k \in \F$ with
	corresponding indices $i_1, \ldots, i_k \in \F$, then they are related to
	the coefficients of the sharing by the $k$ equations
	\begin{align*}
		y_j = \sum_{l=0}^{k-1} c_l i_j^l.
	\end{align*}
	Since this is $k$ linear equations in the $k$ unknowns $c_0, \ldots,
	c_{k-1}$, we know that there is a unique solution. It follows that the same
	is true given our $t$ fixed shares $x$ if we fix a further $k-t$ of them.
	But these additional $k - t$ shares can also be any values, and so we have
	$|\F|^{k-t}$ choices for these, after which all of the shares will be
	fixed. The result follows.
\end{proof}

Next, we present a theorem about the relationship of the shares of a secret to
the coefficients of the associated polynomial. It is clear from the definition
of the shares that they are a linear map of the coefficients, but the following
theorem shows that the reverse is also true; the coefficients can be computed
as a linear map applied to some set of $k$ shares. The proof can be found in
Appendix~\ref{app:shamir}.

\begin{theorem}\label{thm:shareCoeffLO}
	Let $x \in \F$ and ${(x_i)}_{i \in \mathcal{I}}$ be a $k$-sharing of $x$
	for some index set $\mathcal{I} \subset \F$ with $|\mathcal{I}| = n \ge k$.
	Let the associated coefficients for the sharing be $c_0, \ldots, c_{k-1}$.
	Then for each $\mathcal{R} \subset \mathcal{I}$ with $|\mathcal{R}| = k$
	and $j \in \{0, \ldots, k-1\}$, there exist field elements
	${(\lambda_i^{(j)})}_{i \in \mathcal{I}}$ such that
	\begin{align*}
		c_j = \sum_{i \in \mathcal{R}} \lambda_i^{(j)} x_i.
	\end{align*}
\end{theorem}

\subsection{Discrete Logarithm Assumption}

We define here a computational assumption that will be relevant to our random
key pair generation protocol. The assumption is a standard and common one from
the literature: that computing the discrete logarithm is hard. This is
formalised for our specific context as follows.

\begin{assumption}[Discrete Logarithm]\label{ass:dlog}
	Let $p$ be a prime such that $p \ge 2^b$ where $b \in \N$. Let $\G$ be a
	group of prime order $p$ with generator $g$ and let $\F$ be the associated
	finite field of integers modulo $p$. Then for every probabilistic
	polynomial-time Turing machine $T$ and uniformly random field element $x
	\in \F$, $\P(T(\mathbb{G}, g, g^x) = x)$ is negligible.
\end{assumption}

\section{Security}

\newcommand{\exec}[4]{\textup{\texttt{EXEC}}_{#1, #2} (#3, #4)}
\newcommand{\execD}{\exec{\proto}{\mathcal{A}}{x}{z}}

\newcommand{\ideal}[4]{\textup{\texttt{IDEAL}}_{#1, #2} (#3, #4)}
\newcommand{\idealD}{\ideal{\proto}{\mathcal{S}}{x}{z}}

In this section we will outline the security model and definitions that we will
use to prove that the presented protocols are secure. In general, we will
consider an active, rushing adversary with static corruptions. We also define
the adversary to be computationally bounded, i.e.\ so that it cannot compute
discrete logarithms, as in Assumption~\ref{ass:dlog}. In this \paper{} these
properties should be assumed of the adversary unless specified otherwise.

The notion of security we will use is that of Canetti~\cite{c00}, which is
based on the idea of comparing a designed protocol to an ideal case that is
secure by definition. More precisely, we consider an $n$-party function $f$
that takes as input the inputs of each of the parties, and gives as output the
outputs for all of the parties; this defines what functionality we want our
protocol to achieve. Then, we define a protocol $\pi$ which we want to
``securely'' realise $f$, and prove that it is secure by comparing the
``real-life model'' in the presence of an adversary $\mathcal{A}$ and the
``ideal case'':
\begin{itemize}
	\item \textit{Real-life model}: The corrupted parties are controlled by the
		adversary $\mathcal{A}$, which learns their identities and inputs. The
		protocol $\pi$ proceeds in rounds, where in each round uncorrupted
		parties first follow $\pi$ correctly and send any messages they need
		to. Next, $\mathcal{A}$ receives any messages destined for corrupted
		parties, and then decides what messages the corrupted parties should
		send in that round. This process repeats each round until $\pi$ has
		completed execution. The parties then produce their output; uncorrupted
		parties output their true output, while corrupted parties output a
		special symbol $\bot$ to indicate that they were corrupted.

	\item \textit{Ideal case}: All parties hand their inputs for the protocol
		to an incorruptible trusted party $T$ which computes the ideal
		functionality $f$ and then hands each party their respective results.
		Each party produces their respective output as in the real-life model.
\end{itemize}
See the paper for more details~\cite{c00}. An important property to note is
that the adversary gets to see the honest parties' message in a given round
before deciding its own; this property is called \textit{rushing}. We call an
adversary $t$-limited if it controls at most $t$ parties.

If for each $\mathcal{A}$ that operates in the real-life model, we can
construct an adversary $\mathcal{S}$ in the ideal case such that for the random
variable that represents the information gathered by $\mathcal{A}$ and the
outputs of the parties, $\mathcal{S}$ can construct an output with the same
distribution, then we will say that the protocol $\pi$ is secure. Canetti also
proves that this security definition enjoys a \textit{composability} property,
in that if a protocol $\pi$ uses a secure subprotocol $\pi'$ as part of its
execution, then to prove $\pi$ is secure one need only prove this when $\pi'$
has been substituted by the associated $n$-party functionality $f$.
Additionally, in this \paper{} we will consider \textit{active} adversaries,
which means that they can deviate arbitrarily from the protocol and send
arbitrary messages. Before the start of the protocol, the adversary (this
applies to both $\mathcal{A}$ and $\mathcal{S}$) may also modify the inputs of
the corrupted parties in any way.

The random variable corresponding to the execution of the protocol $\pi$ in the
presence of an adversary $\mathcal{A}$ is defined as follows. Denote by
$\texttt{ADVR}_{\pi, \mathcal{A}}(x, z, r)$ the output of $\mathcal{A}$ on
running $\pi$ with auxiliary input $z$ and input $x = (x_1, \ldots, x_n)$ with
random input $r$. This output consists of its auxiliary input, random input,
corrupted parties' inputs, corrupted parties' random inputs, and all messages
sent and received by the corrupted parties during the execution of $\pi$.
Denote by ${\texttt{EXEC}_{\pi, \mathcal{A}}(x, z, r)}_i$ the output of party
$P_i$ at after running $\pi$. Recall that this output will be $\bot$ if $P_i$
is corrupted. Then we have the following definition.

\begin{definition}
	Let $\pi$ be an $n$-party protocol and let $\mathcal{A}$ be a $t$-limited
	adversary. Let $x$ be some input for the parties, and $z$ be some auxiliary
	input. Define $\exec{\pi}{\mathcal{A}}{x}{z}$ to be the joint random variable
	\begin{align*}
		\left(
			\textup{\texttt{ADVR}}_{\pi, \mathcal{A}}(x, z, r),
			{\textup{\texttt{EXEC}}_{\pi, \mathcal{A}}(x, z, r)}_1,
			\ldots,
			{\textup{\texttt{EXEC}}_{\pi, \mathcal{A}}(x, z, r)}_n
		\right),
	\end{align*}
	where $r$ is chosen uniformly randomly.
\end{definition}

The auxiliary input $z$ is used to prove the composability of the security
definition. It represents a possible state for $\mathcal{A}$ that might occur
due to the protocol being a subprotocol of a larger protocol, and so can
include information about previous executions up to that point. Similarly, we
define the random variable that corresponds to the execution in the ideal case
with and adversary $\mathcal{S}$. Denote by $\texttt{ADVR}_{\pi,
\mathcal{S}}(x, z, r)$ the output of $\mathcal{S}$ on running $\pi$ with
auxiliary input $z$ and input $x = (x_1, \ldots, x_n)$ with random input $r$.

\begin{definition}
	Let $f$ be an $n$-party function and let $\mathcal{S}$ be an adversary. Let
	$x$ be some input for the parties, and $z$ be some auxiliary input. Define
	$\ideal{f}{\mathcal{S}}{x}{z}$ to be the random variable
	\begin{align*}
		\left(
			\textup{\texttt{ADVR}}_{\pi, \mathcal{S}}(x, z, r),
			{\textup{\texttt{EXEC}}_{\pi, \mathcal{S}}(x, z, r)}_1,
			\ldots,
			{\textup{\texttt{EXEC}}_{\pi, \mathcal{S}}(x, z, r)}_n
		\right),
	\end{align*}
	where $r$ is chosen uniformly randomly.
\end{definition}

We will often refer to $\exec{\pi}{\mathcal{A}}{x}{z}$ (or sometimes all parts
of it excluding the outputs of the parties) as the \textit{view} of
$\mathcal{A}$ running $\pi$ with inputs $x$ and $z$, as this is what
$\mathcal{A}$ ``sees'' during execution; particularly the messages. Similarly,
we often call the output of $\mathcal{S}$ the \textit{simulated view}.

\begin{definition}\label{def:sec}
	Let $f$ be an $n$-party function and let $\pi$ be an $n$-party protocol. We
	say that $\pi$ $t$-securely evaluates $f$ if for any nonadaptive and
	$t$-limited adversary $\mathcal{A}$ there exists a nonadaptive adversary
	$S$ with running time polynomial in the running time of $\mathcal{A}$ such
	that for all $x$ and $z$
	\begin{align*}
		\exec{\pi}{\mathcal{A}}{x}{z} \eqd \ideal{f}{\mathcal{S}}{x}{z}.
	\end{align*}
\end{definition}

\begin{remark}
	While the auxiliary input $z$ serves an important role in allowing for the
	security definition to be composable, its presence in the individual
	security proofs is not needed; since $\mathcal{S}$ is given $z$, it is
	trivial for it to include it in its output. For this reason, we do not
	mention the auxiliary input in our security proofs.
\end{remark}

\begin{remark}
	The definition presented is a simplified version of the one given in
	Canetti's paper. This is because the latter is made to be general over the
	different possible levels of security: \textit{perfect} security,
	\textit{statistical} security and \textit{computational} security. It is
	also designed to be general over possibly infinite domains for the inputs.
	However, in our case we only consider finite domains and perfect security,
	so the definition is suitably simplified. Note that even though a protocol
	may use a subprotocol that has only statistical or computation security, we
	can still prove it is secure using perfect security, since perfect security
	implies statistical security, which in turn implies computational security.
	For the protocols in this \paper{} we will prove perfect security since it
	is both the strongest result but also more convenient.
\end{remark}

The main difference of this definition of security when compared to earlier
models~\cite{gmr89,mr92,b92} is the combination of the notions of
\textit{secrecy} and \textit{correctness}. The former captures the idea that no
information is gained by the adversary, which corresponds to the distribution
of the messages in the view. The latter captures the fact that the protocol
should correctly do what it is designed to do, and this is captured by having
the outputs of the parties in the view. Thus, Canetti's model captures both of
these properties in the single equality of Definition~\ref{def:sec}, whereas
earlier models required two separate proofs for each of the properties. Note
that this difference is not merely cosmetic; see Canetti's paper~\cite{c00} for
examples and discussion of how a protocol that is intuitively insecure can be
proven to be secure in earlier models but is correctly identified as insecure
by Canetti's security definition.

\section{Cryptographic Primitives}

In this section we briefly outline some cryptographic primitives that will be
used in our SMPC protocols.

\subsection{Public Key Encryption}

\newcommand{\encrypt}[2]{E_{#1}\left(#2\right)}

For some protocols, we will assume the existence of a public key encryption
system. The precise security definition of the encryption is not important for
this context. We will denote the output of encryption of a message $m$ using a
public key $\kappa$ by $\encrypt{\kappa}{m}$.

\subsection{Zero Knowledge Proofs}

We will make use of general purpose zero knowledge (ZK) proofs in some of our
protocols. In some cases, we will require that these proofs have the
\textit{perfect zero knowledge} property, which is summarised in the following
definition by Goldreich~\cite{goldreich_2001}.

\begin{definition}[Perfect Zero Knowledge]
	Let $(P, V)$ be an interactive proof system for some language $L$. We say
	that $(P, V)$ is \textbf{Perfect Zero Knowledge} if for every probabilistic
	polynomial time interactive Turing machine $V^\ast$ there exists a
	probabilistic polynomial time algorithm $M^\ast$ such that for every $x \in
	L$ the following two conditions hold:

	\begin{enumerate}
		\item $M^\ast$ succeeds with probability at least $\frac{1}{2}$; i.e.\
			if we denote a failure by the output of $\bot$, then
			\begin{align*}
				\P(M^\ast(x) = \bot) \le \frac{1}{2}.
			\end{align*}

		\item Let $X$ be the random variable that is the output of $V^\ast$
			after interacting with the interactive machine $P$ on common input
			$x$, and let $m^\ast(x)$ be the random variable that is exactly
			$M^\ast(x)$ but conditioned on the fact that $M^\ast(x) \ne \bot$.
			Then we require that these two random variables are equally
			distributed, i.e.
			\begin{align*}
				m^\ast(x) \eqd X.
			\end{align*}
	\end{enumerate}
\end{definition}

\subsection{Consensus}

\newcommand{\consrng}[2]{\rho_\texttt{RNG}^{#1, #2}}

For some of the protocols described in this \paper{}, we will make use of a
consensus algorithm. A consensus algorithm allows a network of
\textit{processes}, some of which may be \textit{faulty}, to arrive at a joint
decision value. The key properties of a consensus algorithm that we will
consider are those defined Buchman et~al.~\cite{buchman_2018}.

\begin{enumerate}
	\item \textit{Termination}: Every nonfaulty process eventually decides on a
		value.
	\item \textit{Agreement}: No two nonfaulty processes decide on different
		values.
	\item \textit{Validity}: A decided value is valid, i.e., it satisfies the
		predefined predicate denoted $\texttt{valid}()$.
\end{enumerate}

Note that it is not explicitly stated, but it is necessary for the function
$\texttt{valid}$ to be global in the sense that every nonfaulty process
computes the same result for $\texttt{valid}(v)$. This is important to keep in
mind for the following, in which some decision values will contain data
encrypted for specific players, which of course means only those specific
players can decrypt this data and hence the associated plaintext should not be
used as a part of the checks in $\texttt{valid}$.

We define a specific protocol that uses consensus when generating global
(secret shared) random numbers, which we denote $\consrng{n}{k}$. It is defined
as follows. Let the parties participating in the protocol be ${\{P_i\}}_{i \in
\mathcal{I}}$, such that each party $P_i$ has a public and private keypair
where the public key is denoted $\kappa_i$. Each party $P_i$ has input
$c_0^{(i)}, \ldots, c_{t-1}^{(i)}$, which are the coefficients of a polynomial
that is to be used for a threshold $t$ sharing of the number $c_0^{(i)} \in
\F$. They then do the following to construct their input for a consensus
protocol $\rho$:

\begin{enumerate}
	\item Create $n$ shares $r_1^{(i)}, \ldots, r_n^{(i)}$ from the
		coefficients $c_0^{(i)}, \ldots, c_{t-1}^{(i)}$.
	\item Obtain the set of encrypted values $e_i = \{e_1^{(i)},
		\ldots, e_n^{(i)}\}$, where $e_j^{(i)} =
		\encrypt{\kappa_j}{r_j^{(j)}}$.
	\item Create a ZK proof $\zeta_i$ that asserts the following:
		\begin{itemize}
			\item $e_j^{(i)} = \encrypt{\kappa_j}{r_j^{(i)}}$ for all
				$j \in \mathcal{I}$.
			\item The values $r_1^{(i)}, \ldots, r_n^{(i)}$ constitute
				a valid and consistent sharing of some value in $\F$ (in this
				case, that value is $c_0^{(i)}$).
		\end{itemize}
\end{enumerate}

A possible decision value for $\rho$ is a set ${\{(e_i, \zeta_i)\}}_{i \in I}$
where $I \subset \mathcal{I}$. The predicate is $\texttt{valid}_\texttt{RNG}$,
which is true precisely when $|I| > t$ and $\texttt{verify}(\zeta_i) = 1$ for
all $i \in I$. The output for each party is the decrypted set of shares which
were encrypted for their public key; namely the output for party $P_i$ is
${(r_i^{(j)})}_{j \in I}$.

Part of the use of the ZK proofs here is to achieve the goal of having
\textit{verifiable secret sharing}, for which the more common solution is to
use more specialised (and hence usually more efficient) techniques such as that
of Feldman~\cite{feldman_1987} or Pedersen~\cite{p91a}. The reason
that these solutions were not used is to weaken the synchronicity requirements
but also reduce the number of rounds of communication. Using the standard
solutions allows each party to identify when a dealer has not consistently
shared their secret, but each party needs to be aware of this and hence they
need to agree on dealers that are faulty. This usually requires additional
rounds of broadcasting complaints (if there are any), and often doing so using
a secure broadcast channel%
\footnote{%
	In this context, \textit{secure broadcast} means secure in the sense of our
	security definitions, and not in the intuitive encrypted messaging sense; a
	secure broadcast can be thought of as a protocol where each honest party
	outputs the same value. This is obviously not achieved by a na\"{\i}ve
	approach of simply sending the message to everyone, as the adversary may
	send different messages to different parties.
}.
Using the more powerful ZK proofs, in which it is proved that all of the
encrypted shares are consistent, allows each party to check that every other
parties' share is correct in the course of the consensus algorithm. This means
that upon achieving consensus, no further coordinating or complaint
broadcasting is needed as the only decision values that are selected are those
for which the required threshold of parties agreed that $\texttt{valid}$
returned true. Using $\consrng{n}{k}$ thus only introduces the synchronicity
required for the consensus algorithm, which for example in the case of
Tendermint~\cite{buchman_2018} is only partial synchronicity.

This consensus protocol will be used as a subprotocol for some of our SMPC
primitives, and so we will provide some brief justification as to why they can
be used securely. Recall that the two key properties that
definition~\ref{def:sec} captures are secrecy and correctness. The former
requires that the messages do not leak any information, and we can see that
this is the case given that the proposed values only contain encrypted data, ZK
proofs and public keys, all of which do not leak information given appropriate
assumptions. The latter requires that the adversary cannot influence the
output. This is not strictly attained, for example when using
Tendermint~\cite{buchman_2018} consensus a party that is controlled by the
adversary can propose a block and hence influence what is in it, but to counter
this we will just combine the consensus protocol with other protocols to arrive
at adequate security for the larger protocol that uses consensus. Regarding the
computation model, we require that a subprotocol instance can behave like a
``round'' and have a definite completion time. This is ensured for protocols
that want to use consensus by the termination property of the consensus
protocol.

\section{ECDSA Signature Algorithm}

%TODO: Find a reference for the ECDSA algorithm.
The ECDSA algorithm for an elliptic curve is defined as follows. Let an
elliptic curve defined over the field $\mathbb{F}_p$ be given such that it has
a generator $G$ of prime order $n$. Let $d \in \seq{n-1}$ be a private key with
associated public key $dG$, and let $m$ be the message to be signed. Let $z$ be
a hash of $m$ such that the bit length of $z$ is the same as the bit length of
$n$. Then the following steps are used to generate the signature $(r, s)$ of
$m$ under the private key $d$:

\begin{enumerate}
	\item Pick $k$ uniformly randomly from $\seq{n-1}$.
	\item Compute $(x, y) = kG$.
	\item Let $r = x \mod n$.
	\item Compute $s = k^{-1}(z + rd) \mod n$.
	\item Output the signature $(r, s)$.
\end{enumerate}

Note that this is the same standard as in~\cite{gjkr96} except it is adapted
for use with elliptic curves. To implement this signing algorithm in SMPC, the
parties will need to have shares of the private key $d$ such that no small
enough subset knows $d$. This can be achieved by using a random key-pair
generation algorithm. The parties will also need to generate $k$ along with the
associated public key $kG$, which can also be achieved using a random key-pair
generation protocol. Next, the parties will need to perform some arithmetic
operations on the shares of $d$ and $k$ and the hash of the message $z$.
Arithmetic operations on shares can be performed locally when the operation is
addition or when multiplying by a public (not shared) value. This means that
computing $r$ as well as $z + rd$ can be done locally. The two important
operations are thus inverting $k$, and performing a multiplication of the
shared values $k^{-1}$ and $z + rd$.

\subsection{Comparison to~\cite{gjkr96}}

Here we will make an explicit comparison of our approach to~\cite{gjkr96}. The
way that our threshold signature protocol is carried out is the same
as~\cite{gjkr96} when viewed as the composition of subprotocols. That is, if we
consider the subprotocols of random number generation (RNG), random zero
generation (RZG), random key-pair generation (RKPG), multiplication of shares
(including opening the result) and inversion of a shared value, then we use
these subprotocols in the same way as in~\cite{gjkr96} to compute the
signature, with the following small differences:
\begin{itemize}
	\item We target the ECDSA algorithm targeted at using elliptic curves. This
		means that~\cite{gjkr96} instead computes
		\begin{align*}
			r &= \left(g^{k^{-1}} \mod p\right) \mod q\\
			s & = k(m + xr) \mod q,
		\end{align*}
		where $p$, $q$ are primes such that $q \mid p - 1$ and $g$ is an
		element of order $q$ in the multiplicative group of integers modulo
		$p$, $\mathbb{Z}_p^\ast$. This just means that some values live in
		different sets, e.g.\ generating a random key-pair generates a random
		shared value $k \in \seq{q}$ but the public value $g^k$ is in the group
		generated by $g \in \mathbb{Z}_p^\ast$, whereas our public value will
		be a point on the elliptic curve.
	\item Order of subprotocols. When computing a signature, we will need to do
		two multiplications (one is needed during inversion, and then one for
		the final step in computing $s$) and for this we will need two random
		sharings of zero, obtained by calling RZG twice.  In~\cite{gjkr96}
		these two calls to RZG are performed near the start of the protocol, to
		be used later. On the other hand, in our protocol we will include them
		as a part of the multiplication subprotocol. However this has no
		significant affect on the proof (this is true both in~\cite{gjkr96} and
		in this \paper{} where we use Canetti's security framework). We make
		this difference in this \paper{} so that we can prove the security of
		multiplication in Canetti's security framework which allows
		multiplication to be more easily reused in other contexts as a single
		securely composable protocol.
\end{itemize}

The only other important differences between the protocol presented in this
\paper{} and that of~\cite{gjkr96} is how the subprotocols for RNG, RZG, and
RKPG are defined. These subprotocols are specified in Section~\ref{sec:rng},
Section~\ref{sec:rzg} and Section~\ref{sec:rkpg} respectively.

\subsubsection{Security}

One question that arises is the following: if the protocol in this \paper{} has
these differences to the protocol in~\cite{gjkr96}, does the proof of security
presented in~\cite{gjkr96} still apply here? We answer in the affirmative. To
see why, we will consider the differences in turn. First, we turn our attention
to the use in this \paper{} of different realisations of the subprotocols RNG,
RZG and RKPG.\ In this paper, we will prove that our realisation of these
subprotocols are secure under Canetti's definition of security~\cite{c00}.
Notice that the definition of security used by~\cite{gjkr96} is based on older
works~\cite{gmr89,mr92,b92} which provide a weaker definition. Specifically,
Canetti's definition combines both the messages sent by the parties as well as
the outputs of the honest parties as part of the view, whereas the older models
only include the messages in the view (and so require separate proofs for
correctness and secrecy as seen in~\cite{gjkr96}). This means that our security
proofs will imply security in the weaker definition used by~\cite{gjkr96}%
\footnote{%
	Note that in~\cite{gjkr96} there are multiple proofs for different
	adversary models. Here we are specifically interested in the proof
	corresponding to malicious adversaries, which is for protocol
	\texttt{DSS-Tresh-Sig-2} in their paper, as this corresponds to our
	adversary model.
}.
However, in this \paper{} we prove that these subprotocols are secure
\textit{individually} and then make sure of the composability property of
Canetti's security definition to use them together, whereas in~\cite{gjkr96}
the proof of security is given for the \textit{entire} signature protocol,
instead of its constituent parts. We argue that this still means that the proof
in~\cite{gjkr96} is valid if we substitute our subprotocols, because we can
take the simulators used in the proofs of security of $\piRNG{n,k}$,
$\piRZG{n,k}$ and $\piRKPG{n,k}$ and use them to make the same arguments as
their relevant parts in the proof of security in~\cite{gjkr96}. In those places
in the proof where the simulator needs to cheat to cause the output of one of
the subprotocols to be a specific value, we note that this can be achieved with
the simulators in this paper due to the stronger notion of security under which
they are proven.

\section{SMPC Primitives}

In this section we will define the primitives which constitute our SMPC
protocols and prove their security. It is clear that the security model
trivially applies when a protocol is ``local''; i.e.\ does not involve any
exchange of messages. This means that we need only consider proofs of security
for those protocols that involve sending messages. In our case, these protocols
are open (Section~\ref{sec:open}), random number generation
(Section~\ref{sec:rng}), random zero generation (Section~\ref{sec:rzg}), and
random keypair generation (Section~\ref{sec:rkpg}).

\subsection{Open}\label{sec:open}

At some points during a computation on shared values, and almost certainly at
the end of a larger protocol, we will want to reveal the secret corresponding
to the underlying shares. We call this revealing operation ``Opening''.
Probably the most obvious ideal functionality for this would be defined as
follows: $f$ takes as input the shares of the secret from each party and gives
as output the corresponding secret $s$ to each party. However, if we were to
try to achieve this with the simplest and most obvious protocol, everyone
broadcasting their share and then reconstructing, we would not be able to prove
security in the framework we are using because any simulator, knowing only the
corrupted shares and the secret, would not be able to produce the shares of the
honest parties which it would have to as these are messages that are sent
during the protocol. The interpretation of this is clear: this protocol is not
secure because it leaks private information (the input shares) from the honest
parties, or alternatively, because the ideal functionality specifies that this
private information is not revealed. The two solutions for these two
perspectives leads to two different protocols which warrant use in their own
contexts.

First, from the perspective that the simple broadcast protocol should not
reveal the input shares, the solution is to improve the protocol so that this
information isn't revealed. One way that this can be achieved is by first
generating a random $k$-sharing of $0 \in \F$ and adding that to the shares
before broadcasting them; now the messages sent during the protocol are for a
\textit{random} sharing of the secret, as opposed to the specific sharing that
the parties started the protocol with, and a simulator can now generate these
messages with the right distribution knowing only the secret. Second, from the
perspective that the ideal functionality is too restrictive, we can modify it
to also output the shares themselves along with the secret; a simulator for
this ideal functionality can now easily produce an appropriate view for the
simple broadcast protocol.

We can see that in either of the above cases the key difference is whether or
not the original shares are considered safe to reveal. Often they are safe to
reveal, but there are cases where they are not. One example in which they are
not is having $k$-sharings of two private values $a, b \in \F$, and each party
multiplies their shares locally to get a $2k$-sharing of $ab$ and then opening
this value. While it may be ``safe'' to reveal $ab$ (e.g.\ if $a$ is a secret
but $b$ is an unknown uniformly random value), it is well known that this $2k$
sharing is defined by a degree $2k - 1$ polynomial that is not uniformly random
(an easy reason that this is the case is that given it is the product of two
polynomials, it can't possibly be irreducible), and in fact exhibits enough
structure that knowing the shares, one can often easily determine both $a$ and
$b$ individually.

With the above considerations in mind, we will define two protocols for open,
one for each of the cases. In both cases, we will use a Reed-Solomon (RS)
decoding technique, such as Berlekamp-Welch~\cite{welch_1986} or
Gao's~\cite{gao_2003}, to remain fault tolerant in the presence of corrupt
parties broadcasting incorrect shares. For an $(n, k)$ RS code (i.e. $n, k$
Shamir secret sharing) these algorithms are able to detect when there are up to
$d = n - k + 1$ errors and correct up to $\lfloor \frac{d}{2} \rfloor$ errors%
\footnote{%
	Note that these algorithms are also able to locate the errors when they are
	corrected, which in our context means that we would be able to identify
	corrupt nodes.
}.
Notice that RS decoding allows us to recover the underlying polynomial, and so
this means that we not only obtain the secret for our sharing but also all of
the other shares of every other party.

The protocol implementing the basic ideal functionality (that outputs the
secret as well as all of the input shares) $\fOpen{n,k}$ is denoted by
$\piOpen{n,k}$ and is defined as follows.

\begin{enumerate}
	\item Each party $P_i$ broadcasts their input share $x_i$.

	\item Each party performs the reconstruction of the secret using a suitable
		RS decoding algorithm and outputs the result, along with the other
		reconstructed shares.
\end{enumerate}

\begin{theorem}
	Let $t \le \frac{n - k + 1}{2}$. Then the protocol $\piOpen{n,k}$
	$t$-securely evaluates $\fOpen{n,k}$.
\end{theorem}

\begin{proof}
	Since the only messages sent are the input shares and these are included in
	the output of the trusted party, the proof is trivial. The only thing to
	note is that the output needs to include the input shares of all parties,
	including those corrupted parties that may broadcast arbitrary values. This
	is easily overcome by using an RS decoding technique and noting that the
	restriction on $t$ ensures that it will always be possible to reconstruct
	their shares.
\end{proof}

The protocol implementing the stricter ideal functionality $\fSOpen{n,k}$ is
denoted by $\piSOpen{n,k}$ and is the same as the simple version, except before
revealing shares for reconstruction, the sharing is randomised by adding a
random sharing of $0 \in \F$ first. The protocol for generating this
randomising sharing, $\piRZG{n,k}$, is defined in Section~\ref{sec:rzg}. The
protocol $\piSOpen{n,k}$ is defined as follows.

\begin{enumerate}
	\item The parties participate in $\piRZG{n,k}$ to get a share of zero
		$z_i$.

	\item Each party $P_i$ broadcasts their corresponding share $x_i + z_i$.

	\item Each party performs the reconstruction of the secret using a suitable
		RS decoding algorithm and outputs the result.
\end{enumerate}

\begin{remark}\label{rem:openSync}
	We make a note here on synchronicity regarding the use of a RS decoding
	technique. The main feature of the security definition that relates to
	synchronicity is the fact that the protocol proceeds in rounds and each
	round has to be completed before moving on to the next. This would either
	mean waiting for \textit{all} messages from other players, or using a
	timeout after which any messages that have not arrived are given some
	default value to be used instead. However, these alternatives need not be
	used for the final round of each open protocol in which each party sends
	its share to every other party. This is because as soon as a party has at
	least $n - \lfloor \frac{d}{2} \rfloor$ shares, the decoding algorithm will
	be able to successfully recover the secret, setting the yet to be received
	messages as any value. This means that if there are no more than $\lfloor
	\frac{d}{2} \rfloor$ adversaries (and this category also includes
	offline/unresponsive players), then the round is guaranteed to terminate
	according to any termination assumptions of \textit{only the honest
	players}. This means that in practice we will not need to wait for all
	messages (which is not lively) and also will have no need for timeouts.

	Concretely, the following is a way in which the honest parties can be sure
	to terminate without having to have synchronised clocks or use timeouts.
	\begin{enumerate}
		\item Wait until $n - \lfloor \frac{d}{2} \rfloor$ shares have been
			received. This is guaranteed to happen as there are at least this
			many honest parties.
		\item Starting from the previous step and repeating on every new share
			until a reconstruction is possible: set the shares of the parties
			for which no shares have yet been received as $0$, and run the
			reconstruction algorithm.
	\end{enumerate}
	Repeating the last step will eventually terminate with the correct value.
	First, notice that it will never yield an incorrect value. This is because
	no reconstructions are attempted until there are at least $n - \lfloor
	\frac{d}{2} \rfloor$ values, which means at most $\lfloor \frac{d}{2}
	\rfloor$ of them will be from corrupt parties, and therefore in total there
	will be at most $d$ incorrect shares (including the ones that have not yet
	been received). Thus the RS decoding algorithm will always be able to
	detect that there are errors. To see that it will always terminate, we need
	only realise that as soon as we have $n - \lfloor \frac{d}{2} \rfloor$
	correct values the decoding will succeed, and this will eventually happen
	because there are at least this many honest parties.
\end{remark}

\begin{remark}
	When $k = \frac{2n}{3}$, which occurs during multiplication of two
	$k$-sharings when $k = \frac{n}{3}$, we see that this places the security
	requirement at $t < \frac{n}{6}$ (but the protocol will still be safe with
	abort for up to $t < \frac{n}{3}$). We could achieve a more relaxed $t <
	\frac{n}{3}$ requirement if we augment the opening process with ZK proofs
	that the shares a player broadcasts are correct. In this case, as soon as
	$k$ shares with valid proofs have been received, it is safe to use this
	subset to reconstruct the secret and proceed.
\end{remark}

\begin{theorem}\label{thm:so}
	Let $t$ be such that both $t \le \frac{n - k + 1}{2}$ and $t < k - 1$. Then
	the protocol $\piSOpen{n,k}$ $t$-securely evaluates $\fSOpen{n,k}$.
\end{theorem}

\begin{remark}
	Note that we require $t < k - 1$ instead of the usual $k$. This is because
	when we open, every party learns the secret, which we recall is the share
	corresponding to index $0 \in \F$. With this extra share, the adversary can
	actually reconstruct all of the original shares (if we have the worst case
	$t = k - 1$), which we want to avoid.
\end{remark}

\begin{proof}
	\newcommand{\proto}{\piSOpen{n,k}}

	We will give a high level proof here. Full mathematical details can be
	found in Appendix~\ref{app:so}.

	The idea of the proof is as follows. The simulator $\mathcal{S}$ needs to
	construct the view so that it is consistent with the output secret being
	$s$. Since $\piSOpen{n,k}$ randomises the sharing before reconstructing,
	these shares should look random and so $\mathcal{S}$ can simply construct a
	random sharing ${(y_i)}_{i \in \mathcal{I}}$ of $s$ for that part of the
	view. Now, $\mathcal{S}$ needs to make sure that the rest of the view is
	consistent with the real world view: the corrupted parties inputs
	${(x_i)}_{i \in \mathcal{I}_c}$ and random shares of zero ${(z_i)}_{i \in
	\mathcal{I}_c}$ need to satisfy the relationship $x_i + z_i = y_i$ for each
	$i \in \mathcal{I}_c$. Since the input shares are fixed, $\mathcal{S}$
	simply defines $z_i = y_i - x_i$ for each $i \in \mathcal{I}_c$. This will
	have the right distribution because the sharing of $s$ was chosen randomly
	and $t < k - 1$. Define the simulator adversary $\mathcal{S}$ as follows.

	\begin{enumerate}
		\item Let the output from the trusted party be $s$ (by definition it
			will always be the secret corresponding to the sharing $x$).
			$\mathcal{S}$ begins by constructing a random sharing ${(y_i)}_{i
			\in \mathcal{I}}$ of $s$.

		\item $\mathcal{S}$ now defines
			\begin{align}\label{eq:so_z}
				z_i = y_i - x_i \quad \forall i \in \mathcal{I}_c.
			\end{align}

		\item $\mathcal{S}$ now runs $\mathcal{A}$ after giving it ${(z_i)}_{i
			\in \mathcal{I}_h}$ and ${(y_i)}_{i \in \mathcal{I}_h}$ to get the
			messages ${(m_i)}_{i \in \mathcal{I}_c}$ that $\mathcal{A}$ sends
			in the broadcast round.

		\item $\mathcal{S}$ then outputs
			\begin{align*}
				{(z_i)}_{i \in \mathcal{I}_c},
				{(m_i)}_{i \in \mathcal{I}_c},
				{(y_i)}_{i \in \mathcal{I}_h}
			\end{align*}
			as the generated view.
	\end{enumerate}

	We will now justify why the view of $\mathcal{S}$, $\idealD$, has the
	correct distribution. First, in the real world view the shares of zero
	${(z_i)}_{i \in \mathcal{I}_c}$ will look completely random as there are
	less than $k - 1$ of them, whereas they are shared with threshold $k - 1$.
	In the simulated view this is also true, based on their definition in
	Eq.~\eqref{eq:so_z} and the fact that ${(y_i)}_{i \in \mathcal{I}_c}$ are
	also completely random by construction. Next, the messages from the
	adversary ${(m_i)}_{i \in \mathcal{I}_c}$ will have the correct
	distribution because they were obtained by directly running $\mathcal{A}$
	and all of the information that $\mathcal{A}$ has obtained up until this
	point, namely ${(z_i)}_{i \in \mathcal{I}_c}$, has the right distribution.
	Finally, the last element has the correct distribution because of its
	construction and the relationship Eq.~\eqref{eq:so_z}.

	As for correctness of the output, this is an easy consequence of the
	restrictions on $t$. These restrictions ensure that the honest parties can
	correctly reconstruct the secret using RS decoding.
\end{proof}

\subsubsection{Directed Open}

The protocols $\piOpen{n,k}$ and $\piSOpen{n,k}$ reveal the shared secret to
all players. However, sometimes we will only want to reveal the secret to a
specific player. In this case we use a \textit{directed open} protocol. The
only difference in this case is that instead of sending the shares to everyone,
each player sends their share to the specified player only. We denote the
directed version of $\piOpen{n,k}$ as $\piDOpen{n,k,i}$ and the directed
version of $\piSOpen{n,k}$ as $\piSDOpen{n,k,i}$, where in each case $i \in
\mathcal{I}$ is the index of the player which outputs the secret. The ideal
functionalities $\fDOpen{n,k,i}$ and $\fSDOpen{n,k,i}$ are also defined
correspondingly.

The security of these directed open protocols is summarised in the following
two theorems; they enjoy the same security as their undirected counterparts.
The proofs are almost identical and so are omitted here.

\begin{theorem}
	Let $t \le \frac{n - k + 1}{2}$. Then the protocol $\piDOpen{n,k}$
	$t$-securely evaluates $\fDOpen{n,k}$.
\end{theorem}

\begin{theorem}
	Let $t$ be such that both $t \le \frac{n - k + 1}{2}$ and $t < k - 1$. Then
	the protocol $\piSDOpen{n,k}$ $t$-securely evaluates $\fSDOpen{n,k}$.
\end{theorem}

\subsection{Random Number Generation}\label{sec:rng}

Generating shares of a uniformly distributed (but unknown to all parties)
random number is key for many SMPC protocols. In this section we will describe
how we achieve this and prove the security of our protocol. Security requires
that the output cannot be influenced, and to construct a protocol that meets
this requirement we will first create a protocol that does not meet all of the
requirements for our security framework. We will then utilise this in creating
a fully secure protocol.

\subsubsection{Biased RNG}

The goal of random number generation is for each of the parties to end up with
shares of a global random number that none of the parties know. This technique
was first introduced by Pedersen~\cite{p91a} and works as follows. Each player
begins by generating a random number and creating shares of it. These shares
are then sent too all other players, after which each player should have a
share of the random numbers of the other players, as well as one of their own.
These shares are then summed to produce a single share of the sum of each
players random number. If at least one of the players that participated picked
their random number uniformly randomly, then the global random number will also
be uniformly random.

The main challenge of using this protocol in the presence of malicious
adversaries is that of having the honest parties select the same subset of
shares to use in the sum; not doing so would lead to the honest parties holing
an inconsistent sharing. The way that a corrupted node can cause this to happen
is simply sending their shares to a subset of the parties, and not sending any
to the others. One possible solution to this is to have parties also send
acknowledgement messages upon receiving correct shares on a secure broadcast
channel~\cite{gjkr96}. However, since we are not using a secure broadcast
channel we use an alternative solution: a consensus protocol. The idea is that
once consensus is achieved, all of the honest parties will agree on which
subset of shares from other parties to include in the sum, and will therefore
end up with consistent shares of the same global random number. The protocol is
denoted $\piBRNG{n,k}$ proceeds as follows:

\begin{enumerate}
	\item For all $i \in \mathcal{I}$, player $P_i$ picks $r_i \in \F$
		uniformly randomly and creates shares ${(r_{i,j})}_{j \in
		\mathcal{I}}$.  Denote the coefficients of the random polynomial that
		determines this sharing by $c_0^{(i)}, \ldots, c_{t-1}^{(i)}$, so that
		$r_i = c_0^{(i)}$ and
		\begin{align*}
			r_{i, j} = \sum_{k=0}^{t-1} c_k^{(i)} j^k.
		\end{align*}

	\item Each player $P_i$ then participates in $\consrng{n}{k}$ using the
		generated coefficients, obtaining the output ${(r_{j, i})}_{j \in
		I}$ where $I \subset \mathcal{I}$ is defined implicitly by
		$\consrng{n}{k}$ as the set of players whose shares were agreed to be
		used by $\consrng{n}{k}$.

	\item The final output of player $P_i$ is
		\begin{align*}
			y_i = \sum_{j \in I} r_{j, i}.
		\end{align*}
\end{enumerate}

If we attempt to prove that this protocol (or similar ones such
as~\cite{gjkr96}) Canetti's security model, we encounter a problem. This
problem arises when we define the most obvious ideal functionality for
$\piBRNG{n,k}$. This obvious ideal functionality $\fRNG{n,k}$ is: each party
gives no input and the output is a uniformly random sharing of a uniformly
random number, the output for each party being their respective share. The key
here is that we have specified that not only is the global random number
uniformly random, but also the \textit{sharing} of this random number is
random; each possible sharing of the random number is equally as likely.
Unfortunately, this latter condition is easily thwarted by a rushing adversary,
as they can observe the honest parties' shares first, and then pick their
random numbers associated sharings to be such that their final shares (the sum
of the shares from the chosen subset of parties) are biased. For example, when
using Tendermint for the consensus protocol, the proposer of a block may be a
corrupted party. If this is the case, they are in fact able to realise the
rushing property and can pick their own sharing to be such that they end up
with any share of their choosing. This clearly violates the property that the
sharing of the random number is itself random. In order to overcome this
technical difficulty, we will not prove that $\piBRNG{n,k}$ is secure in
Canetti's security framework but rather just prove its key properties that will
be relevant when we use $\piBRNG{n,k}$ to construct an unbiased protocol
$\piRNG{n,k}$ that will capable of realising the ideal functionality
$\fRNG{n,k}$.

We are interested in the following properties.

\begin{itemize}
	\item \textit{Agreement}: All honest parties output consistent shares of a
		field element.

	\item \textit{Global randomness}: The secret corresponding to the output
		shares is uniformly randomly distributed.

	\item \textit{Secrecy}: No subset of less than $k$ players knows anything
		about the shared secret other than the fact that it is uniformly
		randomly distributed. Further, the shares of the honest parties are
		uniformly random, constrained only by the fact that they are consistent
		with the shares of the corrupted parties.
\end{itemize}

\begin{theorem}
	Let $\mathcal{A}$ be a $t$-limited adversary where $t < k$. Then protocol
	$\piBRNG{n,k}$ satisfies agreement, global randomness and secrecy.
\end{theorem}

\begin{proof}
	We will prove each property in turn.

	\begin{itemize}
		\item \textit{Agreement}: The fact that all honest parties will agree
			on the same subset $I \subset \mathcal{I}$ of parties whose shares
			are to be included in the final sum is guaranteed by the
			termination and agreement properties of the consensus protocol. The
			consensus protocol also ensures (by dint of the validity property)
			that for each $i \in I$ all honest parties will receive consistent
			shares of a field element chosen by $P_i$. It follows that when
			each honest party obtains its output by locally summing each of the
			shares that came from parties in $I$, this sum will also be part of
			a consistent sharing of the sum of the secret random numbers.

		\item \textit{Global randomness}: Since $t < k$ and $k$ sharings are
			used in the final sum, at least one of the secrets used in the sum
			will have been generated by an honest party and hence be uniformly
			random. The property then follows from Corollary~\ref{cor:unif}.

		\item \textit{Secrecy}: We will begin by showing the second property
			required for secrecy: that the shares of the honest parties appear
			uniformly random, given the constraints. The first property will
			follow as a consequence. Let
			$y_c = {(y_i)}_{i \in \mathcal{I}_c}$
			and
			$y_h = {(y_i)}_{i \in \mathcal{I}_h}$,
			and let
			$x = {(x_i)}_{i \in \mathcal{I}_h}
			\in
			\shareRestrict{y_c}{\mathcal{I}}{\mathcal{I}_c}$
			be arbitrary. Let $m \in I$ be such that $P_m$ is honest. This is
			always possible because $|I| > t$. Let
			$K \subset \mathcal{I}_h$
			be such that $|K| = k - t$. We will need to use the fact that
			${(r_{m, i})}_{i \in K}$
			is independent from
			${(r_{j, i})}_{i \in K}$
			for all $j \in I \setminus \{m\}$. To see this, realise that
			${(r_{m, i})}_{i \in \mathcal{I}_c \cup K}$
			is uniformly randomly distributed in $\F^k$, which follows
			from Theorem~\ref{thm:shareDist}. Then the independence follows
			from this and the fact that the adversary chooses its shares
			knowing only the subset
			${(r_{m, i})}_{i \in \mathcal{I}_c}$
			of the shares from $P_m$.

			Now, from Theorem~\ref{thm:restrictSize} we know that
			$|\shareRestrict{y_c}{\mathcal{I}}{\mathcal{I}_c}| = |\F|^{k-t}$.
			Thus, we want to show that
			\begin{align*}
				\P(x \mid y_c) = |\F|^{t-k}.
			\end{align*}
			But again, we know that $x$ will be completely determined by $y_c$
			and any $k-t$ of the shares in $x$, so we can write
			\begin{align}\label{eq:brngSecrecy}
				\P(x \mid y_c) = \P({(x_i)}_{i \in K} \mid y_c).
			\end{align}
			Since ${(r_{m, i})}_{i \in K}$ is independent from ${(r_{j, i})}_{i
			\in K}$ for all $j \in I \setminus \{m\}$, it follows that we can
			apply Corollary~\ref{cor:unif} to see that ${(y_i)}_{i \in K}$ is
			uniformly distributed in $\F^{k-t}$. This fact combined with
			Eq.~\eqref{eq:brngSecrecy} gives the desired result.

			The first property required for secrecy now follows easily: we saw
			above that ${(y_i)}_{i \in K}$ is uniformly distributed, and hence
			for any particular $i \in K$ it follows that $y_i$ is uniformly
			distributed. However, we would theoretically have $0 \in K$, in
			which case $y_0$ would be the secret for the sharing, and hence
			this too is uniformly random conditioned on $y_c$.
	\end{itemize}
\end{proof}

\subsubsection{Unbiased RNG}

Unbiased random number generation is a protocol $\piRNG{n,k}$ that generates a
global random number that is unknown to all parties, but for which the parties
hold shares. It can be described by the $n$-party function $\fRNG{n,k}$ which
takes no input from the parties, and gives output that is $n$ shares of a
uniformly random sharing of a uniformly random number $r$ (unknown to all
parties) with reconstruction threshold $k$, giving each share to the
corresponding party.

To create the unbiased protocol $\piRNG{n,k}$ from the biased protocol
$\piBRNG{n,k}$ the idea is simple. Run $\piBRNG{n,k}$ $k$ times, so that each
of the random secrets represents a coefficient of a uniformly random polynomial
$p$. Then, compute $P_i$'s share of $p(i)$ for each $i \in \mathcal{I}$ by a
linear combination of the shares of the $k$ coefficients, and open these
resulting shares of shares to each corresponding party using a directed open.
Since creating the shares of $p(i)$ is a linear combination of the
coefficients, this can be computed on the shares locally. The protocol is as
follows.

\begin{enumerate}
	\item\label{enum:rngCoeff} The players invoke $\piBRNG{n,k}$ $k$ times. Let
		the $k$ output sharings be denoted by ${\left(r_i^{(1)}\right)}_{i \in
		\mathcal{I}}, \ldots, {\left(r_i^{(k)}\right)}_{i \in \mathcal{I}}$,
		where player $P_i$ receives the shares $r_i^{(j)}$ for all $1 \le j \le
		k$. These represent shares of the coefficients $c_0, \ldots, c_{k-1}$
		of a random degree $k-1$ polynomial $p$.

	\item\label{enum:rngrij} Each party $P_i$ locally computes
		\begin{align*}
			r_{i,j} = \sum_{l=0}^{k-1} r_i^{(l)} j^l
		\end{align*}
		for all $j \in \mathcal{I}$. This defines the sharings ${(r_{i,j})}_{i
		\in \mathcal{I}}$ for each $j \in \mathcal{I}$ which are shares of
		$r_j$ where
		\begin{align*}
			r_j = \sum_{l=0}^{k-1} c_l j^l.
		\end{align*}
		That is, $r_j = p(j)$, or in other words the share for party $P_j$
		corresponding to $p$.

	\item\label{enum:rngDO} The players invoke $\piDOpen{n,k,j}$ $n$ times: for
		each $j \in \mathcal{I}$, $r_j$ is opened towards party $P_j$. Each
		party $P_j$ then finishes by outputting $r_j$.
\end{enumerate}

\begin{theorem}\label{thm:rng}
	Let $k > t$. Then the protocol $\piRNG{n,k}$ $t$-securely evaluates
	$\fRNG{n,k}$.
\end{theorem}

\begin{proof}
	\newcommand{\proto}{\piRNG{n,k}}

	We will present a high level proof here. Full mathematical details can be
	found in Appendix~\ref{app:rng}.

	The idea behind the simulator is as follows. The simulator can run the
	protocol as usual up until step~\ref{enum:rngDO}. At this point,
	$\mathcal{A}$ does not know what the secrets for the random sharings are,
	and the honest parties shares will be uniformly random conditioned on the
	fact that they are consistent with $\mathcal{A}$'s shares. This allows the
	simulator to simply extend the corrupted parties' shares (which are known
	by the simulator since it holds enough shares to reconstruct them) to valid
	random sharings of the target output shares, and this will have the correct
	distribution.

	The simulator is defined as follows.
	\begin{enumerate}
		\item The trusted party is invoked first. Let the output shares of the
			trusted party be ${(y_i)}_{i \in \mathcal{I}}$.

		\item $\mathcal{S}$ runs $\mathcal{A}$ and acts on behalf of the honest
			parties for Step~\ref{enum:rngCoeff}. In doing so, $\mathcal{S}$
			learns the biased shares ${\left(r_i^{(j)}\right)}_{j \in
			\seqZ{k-1}}$ for each $i \in \mathcal{I}_c$, i.e.\ for each
			corrupted party. Label the messages sent and received by
			$\mathcal{A}$ during the invocations of $\piBRNG{n,k}$ as $m$, and
			the messages that $\mathcal{A}$ would send to the invocations of
			$\piDOpen{n,k,j}$ as $m_o$.

		\item $\mathcal{S}$ constructs the honest parties' shares
			${\left({\left(r_i^{(j)}\right)}_{i \in \mathcal{I}_h}\right)}_{j
			\in \seqZ{k-1}}$ as follows. First, $\mathcal{S}$ picks random
			coefficients ${(c_i)}_{i \in \seqZ{k-1}}$ that are consistent with
			the target shares ${(y_i)}_{i \in \mathcal{I}_c}$, i.e.\
			$\mathcal{S}$ picks uniformly randomly from the set
			\begin{align*}
				\set{%
					{(c_i)}_{i \in \seqZ{k-1}} \in \F^k
				}{%
					y_i = \sum_{j=0}^{k-1} c_j i^j
					\;
					\forall i \in \mathcal{I}_c
				}.
			\end{align*}
			Next, the honest parties' shares are chosen at random under the
			condition that ${\left(r_i^{(j)}\right)}_{i \in \mathcal{I}}$ is a
			consistent sharing of $c_j$ for each $j \in \seqZ{k-1}$. Precisely,
			the shares are chosen uniformly randomly from the set
			\begin{align*}
				\set{%
					{\left(
						{\left(
							r_i^{(j)}
						\right)}_{i \in \mathcal{I}_h}\right)
					}_{j \in \seqZ{k-1}}
					\in
					\F^{k(n-t)}
				}{%
					\begin{matrix}
						\forall j \in \seqZ{k-1}\\
						\exists {(c_l')}_{l \in \seq{k-1}}\\
						\forall i \in \mathcal{I}
					\end{matrix}
					:
					r_i^{(j)} = c_j + \sum_{l=1}^{k-1} c_l' i^l
				}.
			\end{align*}
	\end{enumerate}

	We will justify why $\mathcal{S}$ produces a view with the same
	distribution as in the real world case. The parts of the view that consist
	of messages that the corrupted party sends should have the correct
	distribution as they are generated by running a copy of $\mathcal{A}$
	directly, and by ensuring that the preceding information that $\mathcal{A}$
	gains has the right distribution. Also, the messages that the corrupted
	parties receive from the invocations of $\piBRNG{n,k}$ will clearly have
	the right distribution as they are generated by $\mathcal{S}$ invoking
	$\piBRNG{n,k}$ directly, acting on behalf of the honest players. Thus the
	two important parts of the view that we need to make sure have the right
	distribution are the messages the corrupted parties receive during the
	directed opens, and the output shares. Since the latter is determined by
	the former, we need only consider the former (messages received during the
	directed opens). In the real world view $\execD$, these are just uniformly
	random shares of uniformly random numbers conditioned on the fact that they
	be consistent with the corrupted parties shares. This follows from the fact
	that this is true for the honest parties' shares from $\piBRNG{n,k}$ (from
	the secrecy property), and that the shares in the directed opens are just
	linear combinations of these shares. In the ideal case view $\idealD$,
	these shares are constructed by working backwards from the output shares.
	Since the output sharing is uniformly random, it follows that the
	coefficients ${(c_i)}_{i \in \seqZ{k-1}}$ that $\mathcal{S}$ constructs
	(which correspond the random numbers for the $k$ invocations of
	$\piBRNG{n,k}$) will also be uniformly random. When $\mathcal{S}$ extends
	the corrupted parties' shares to open to these values, it does this by
	picking uniformly randomly from the possible sharings, and this ensures
	that the distribution will be the same as in the real world view.
\end{proof}

\subsection{Random Zero Generation}\label{sec:rzg}

Random zero generation is almost the same as random number generation as in
Section~\ref{sec:rng}, except instead of obtaining shares of a global random
number, the parties obtain a sharing of $0 \in \F$. Specifically, the ideal
functionality $\fRZG{n,k}$ takes no inputs and outputs to each party $P_i$ a
share $z_i$ of $0 \in \F$, such that this sharing is uniformly randomly
distributed. The protocol is also almost the same as $\piRNG{n,k}$, except that
instead of generating $k$ random shared coefficients, we only generate $k-1$
because our secret (constant term in the polynomial) is fixed and not random.
This is made precise in the following protocol $\piRZG{n,k}$:

\begin{enumerate}
	\item The players invoke $\piBRNG{n,k}$ with no input $k-1$ times. Let the
		$k-1$ output sharings be denoted by ${\left(r_i^{(1)}\right)}_{i \in
		\mathcal{I}}, \ldots, {\left(r_i^{(k-1)}\right)}_{i \in \mathcal{I}}$,
		where player $P_i$ receives the shares $r_i^{(j)}$ for all $1 \le j \le
		k-1$. These represent shares of the coefficients $c_1, \ldots, c_{k-1}$
		of a random degree $k-1$ polynomial with a zero constant term.

	\item Each party $P_i$ locally computes
		\begin{align*}
			r_{i,j} = \sum_{l=1}^{k-1} r_i^{(l)} j^l
		\end{align*}
		for all $j \in \mathcal{I}$. This defines the sharings ${(r_{i,j})}_{i
		\in \mathcal{I}}$ for each $j \in \mathcal{I}$ which are shares of
		$r_j$ where
		\begin{align*}
			r_j = \sum_{l=1}^{k-1} c_l j^l.
		\end{align*}
		That is, $r_j$ is the share for party $P_j$ corresponding to the
		polynomial with coefficients $c_1, \ldots, c_{k-1}$.

	\item The players invoke directed open $n$ times: for each $j \in
		\mathcal{I}$, $r_j$ is opened towards party $P_j$. Each party $P_j$
		then finishes by outputting $r_j$.
\end{enumerate}

The security theorem and proof for $\piRZG{n,k}$ is nearly identical to that
for $\piRNG{n,k}$, and so is omitted.

\begin{theorem}
	Let $t < k - 1$. Then the protocol $\piRZG{n,k}$ $t$-securely evaluates
	$\fRZG{n,k}$.
\end{theorem}

\begin{remark}
	Note that the requirement on $t$ is that it is less than $k - 1$, as
	opposed to $k$ as was the case for $\piRNG{n,k}$. This is because we are
	generating a sharing of a known, fixed field element which gives one extra
	share of information to the adversary (recall that the secret of a sharing
	is equal to the share corresponding to index $0 \in \F$).
\end{remark}

\subsection{Random Keypair Generation}\label{sec:rkpg}

Random key pair generation is again very similar to random number generation
defined in Section~\ref{sec:rng}, except instead of only obtaining shares of a
global random number, the parties also obtain the public key that corresponds
to the shared random number (private key). Specifically, the ideal
functionality $\fRKPG{n,k}$ takes no inputs and outputs to each party $P_i$ a
share $x_i$ of some uniformly random $x \in \F$ and $y$ such that $y = g^x \in
\G$ where $\G$ is a group with generator $g$.  Additionally, to enable the
proof of security to work each party will also output $g^{x_i}$ for each $i \in
\mathcal{I}$. This modification will be discussed after presenting the
passively secure protocol ${\piRKPG{n,k}}'$ for $\fRKPG{n,k}$. We begin with
this passively secure protocol and then discuss how it can be augmented to
achieve active security later. The protocol is defined as follows.

\begin{enumerate}
	\item The parties invoke $\piRNG{n,k}$, so that party $P_i$ receives the
		share $x_i$.

	\item\label{enum:rkpgPassBC} Each party $P_i$ sends $g^{x_i}$ to every
		other party.

	\item Each party then reconstructs $g^x$ ``in the exponent'', along with
		$g^{x_i}$ for each $i \in \mathcal{I}$. These, along with the share
		$x_i$ constitute the output of the protocol for player $P_i$.
\end{enumerate}

The reason that the public values $g^{x_i}$ for each $i \in \mathcal{I}$ are
included in the output for each party is to allow the simulator to construct a
consistent view. If the simulator did not know these values, then under
Assumption~\ref{ass:dlog} it would have no hope of constructing $g_i \in \G$
such that $g_i = g^{x_i}$ for each $i \in \mathcal{I}_h$. However, for the same
reason the discrete logarithm problem also implies that learning each $g_i$
should not be a concern, and so we are happy to include them in the output of
the protocol. We will formalise this reasoning somewhat after presenting the
security theorem.

\begin{theorem}\label{thm:rkpg}
	Let $k > t$. Then the protocol ${\piRKPG{n,k}}'$ $t$-securely evaluates
	$\fRKPG{n,k}$ in the presence of a passive adversary.
\end{theorem}

\begin{proof}
	The proof is straightforward; since everything in the view of the adversary
	is contained in the output of the trusted party, and since the adversary is
	passive, constructing a simulated view from this trusted output is trivial.
\end{proof}

We now argue that for a $k$-sharing of $x \in \F$, knowing a subset of less
than $k$ of these shares and $g^{x_i}$ for all shares $x_i$, as well as knowing
$g^x$, it is not possible to discover $x$ if the discrete logarithm problem is
hard. This is summarised in the following theorem.

\begin{theorem}
	Let $\G$ be a group of size at least $2^b$ with generator $g$ and prime
	order, such that $\F$ is the finite field associated with this prime. Let
	$x \in \F$ be uniformly random, and $x_1, \ldots, x_n$ be a uniformly
	random $k$-sharing of $x$ where $k, n \in \N$ and $k \le n$. Let $t \in \N$
	and $I \subset \seq{n}$ be given such that $t < k $ and $|I| = t$. Suppose
	that a polynomial time Turing machine $\mathcal{D}$ takes as input $\G$,
	$g$, ${(x_i)}_{i \in I}$, ${(g^{x_i})}_{i \in \seq{n}}$ and $g^x$, and
	outputs $z \in \F$. Then if Assumption~\ref{ass:dlog} holds, it follows
	that for all such $\mathcal{D}$, $\P(z = x)$ is negligible in $b$.
\end{theorem}

\begin{proof}
	As is standard for these kinds of results, we will proceed by
	contradiction; assume that there exists $\mathcal{D}$ that outputs $z$ such
	that $\P(z = x)$ is non-negligible in $b$. We will show how this can be
	used to solve a target instance of the discrete logarithm problem. To this
	end, let $y \in \G$ be arbitrary. Pick $t$ independent and uniformly random
	elements of $\F$ $x_1, \ldots, x_t$, and $k - t - 1$ independent and
	uniformly random elements of $\G$, labelled $g_{t+1}, \ldots, g_{k-1}$.
	Similarly label $g_i = g^{x_i}$ for each $i \in \seq{t}$. Also extend the
	labels of $x_i$ so that for $i \in \{t+1, \ldots, k-1\}$ $x_i$ is the
	unique value such that we have $g_i = g^{x_i}$ and also define $x_0$ to be
	the unique value that satisfies $y = g^{x_0}$. Note that these values are
	well defined but not known by any polynomial time Turing machine, they are
	merely defined for convenience of exposition. From
	Theorem~\ref{thm:shareCoeffLO} we know that there exist values
	$\lambda_0^{(i)}, \ldots, \lambda_{k-1}^{(i)}$ such that
	\begin{align*}
		c_i = \sum_{j=0}^{k-1} \lambda_j^{(i)} x_j
		\quad
		\forall i \in \seqZ{k-1},
	\end{align*}
	where $c_0, \ldots, c_{k-1}$ are the uniquely defined coefficients that
	correspond to a $k$-sharing that is consistent with the shares $x_0,
	\ldots, x_{k-1}$. Using these values we can compute
	\begin{align*}
		g^{c_i} = h_i = \prod_{j=0}^{k-1} {g_j}^{\lambda_j^{(i)}}
	\end{align*}
	for all $i \in \seqZ{k-1}$. This in turn allows us to compute the remaining
	shares in the exponents; we can compute
	\begin{align*}
		g_i = \prod_{j=0}^{k-1} {h_j}^{i^j}
		\quad
		\forall i \in \{k, \ldots, n\}.
	\end{align*}
	By construction, the shares corresponding to $g_1, \ldots, g_n$ are a
	consistent sharing of $x_0$ and agree with the subset $x_1, \ldots, x_t$.
	We may now run $\mathcal{D}$ with input $\G$, $g$, ${(x_i)}_{i \in
	\seq{t}}$, ${(g_i)}_{i \in \seq{n}}$, and $y$. The output is $z$ where
	$\P(z = x_0)$ with non-negligible probability. Since $y = g^{x_0}$, this
	completes the proof.
\end{proof}

We now discuss how to augment ${\piRKPG{n,k}}'$ to be secure in the presence of
an active adversary. This is not too difficult to achieve; the only part of the
passively secure protocol in which the adversary can send (modified) messages
is in step~\ref{enum:rkpgPassBC}. This can hinder the reconstruction of the
global public key and also the public keys corresponding to the shares if there
is no way to detect that these values are incorrect. Thus, we need only include
a way to ensure that honest parties can detect correct values, and after
receiving $k$ correct values they will be able to perform the reconstruction.
To achieve this required detection, we will augment $\piRNG{n,k}$ to output,
along side the usual random shares, a perfectly hiding commitment to each of
the shares of the parties. Then, ${\piRKPG{n,k}}'$ is augmented by also
submitting in step~\ref{enum:rkpgPassBC} a zero knowledge proof that has the
perfect zero knowledge property that the value they sent corresponds to the
commitment to their share. This will enable honest parties to detect which
broadcasted values are correct. The associated simulator for the proof would
then be able to work as previously, except to generate the ZK proof messages it
can leverage the perfect zero knowledge property. The augmented protocol
$\piRKPG{n,k}$ is as follows.

\begin{enumerate}
	\item The parties invoke the augmented RNG protocol ${\piRNG{n,k}}'$, so
		that party $P_i$ receives the share $x_i$. Every party also receives
		for each $i \in \mathcal{I}$ the commitment $c_i$ to $x_i$.

	\item Each party $P_i$ sends $g^{x_i}$ to every other party, along with a
		ZK proof $z_i$ that attests to the fact that the discrete log of
		$g^{x_i}$ is the same as the value committed to by $c_i$.

	\item Each party then reconstructs $g^x$ ``in the exponent'', along with
		$g^{x_i}$ for each $i \in \mathcal{I}$. These, along with the share
		$x_i$ constitute the output of the protocol for player $P_i$.
\end{enumerate}

\subsection{Local Arithmetic}

We will now describe the fundamental computational primitives for the SMPC
algorithm. This includes the standard operations that can be performed on
elements of a field, i.e.\ the field operations and their inverses. For most of
the operations, applying them locally to the shares (that is, each party
applies them to their own shares, and on a collective level we can describe
this as operating element-wise) results in the equivalent effect on the secret
itself, with no effect to the threshold of the sharing. The one operation for
which this is not the case is multiplication; here operating locally will
indeed give a share of the product of the secrets, but the threshold will
increase and the sharing will exhibit unwanted structure. This means some extra
work will need to be done to achieve a desired multiplication protocol. For all
of the other operations, the fact that they can be carried out locally means
that no messages are exchanged between parties and therefore a ``protocol''
that encapsulates one of these operations is trivially secure under the
security definition, and hence the corresponding theorems will not be stated
nor proven.

To simplify notation, write the sharing ${(x_i)}_{i \in \mathcal{I}}$ as
$x_\mathcal{I}$. We will denote local operations on shares as follows, where on
the left we have the notation and the right we have what it means:
\begin{align*}
	c_\mathcal{I} &= a_\mathcal{I} + b_\mathcal{I} &
	c_\mathcal{I} &= {(a_i + b_i)}_{i \in \mathcal{I}}\\
	c_\mathcal{I} &= a_\mathcal{I} - b_\mathcal{I} &
	c_\mathcal{I} &= {(a_i - b_i)}_{i \in \mathcal{I}}\\
	c_\mathcal{I} &= a_\mathcal{I} b_\mathcal{I} &
	c_\mathcal{I} &= {(a_i b_i)}_{i \in \mathcal{I}}\\
	c_\mathcal{I} &= -a_\mathcal{I} &
	c_\mathcal{I} &= {(-a_i)}_{i \in \mathcal{I}}\\
	c_\mathcal{I} &= r a_\mathcal{I} &
	c_\mathcal{I} &= {(r a_i)}_{i \in \mathcal{I}}
\end{align*}
We will consider an example to make this clear. Consider addition of two shared
values where each party $P_i$ holds the shares $a_i$ and $b_i$. Then $P_i$
simply defines the share $c_i$ to be $a_i + b_i$, and this will be its output
of the addition protocol.

\section{SMPC Protocols}

In this section we will outline the protocols that are built from the
primitives outlined in the preceding section. The main goal that we are working
towards is to be able to perform an ECDSA signature where the private key is
distributed among the parties as a shared secret. To do this, we will construct
the SMPC protocols that will be sufficient to perform the signature. We will
then be able to leverage the composability property of the security framework
to be able to compose our building blocks together in a secure way. However,
this security definition does not capture all of the security properties that
we seek in general. For instance, it does not capture the security of our
secret (in the context of Shamir Secret Sharing). We can see this by
considering the following obviously insecure protocol for multiplying two
secrets and then immediately opening the result: simply open the two sharings
and then multiply them together locally. The reason that this continues to
satisfy security as defined in Definition~\ref{def:sec}, despite being
intuitively insecure, is that it is evaluating an ideal functionality that we
do not actually want (the ideal functionality that we want simply takes the
input shares and gives the product of the secrets). Rather, it evaluates
precisely the secure functionality defined by this example protocol: take the
input shares and give the product of the secrets \textit{but also the input
secrets themselves} (the latter because of the opens that are performed in the
protocol description). Now, it becomes clear that the reason this example
protocol is secure under Definition~\ref{def:sec}, but not intuitively, is
because the actual protocols explicitly describes itself as revealing
information. We can see that the problem in this case arises because of opening
values that should not be revealed. The solution to this is to make sure that
whenever we perform an open during our protocol, the revealed value should be
uniformly randomly distributed, so that nothing is ``learned'' from it. Because
open is the only protocol that reveals information about secrets, as long as we
make sure that we do this when constructing a protocol out of our primitive
protocols, we can be sure that we are securely evaluating what we intend
without worrying about exactly how we got there.

\subsection{Multiply and Open}

The ideal functionality $\fMulOpen{n,k}$ that represents the multiply and then
open protocol takes as input two sets of shares for two field elements and
gives as output to each party the product of these two secrets.

Let $a_\mathcal{I}$ and $b_\mathcal{I}$ be two $k$-sharings of respective field
elements $a, b \in \F$. The protocol $\piMulOpen{n,k}(a_\mathcal{I},
b_\mathcal{I})$ is defined by
\begin{align*}
	c_\mathcal{I} &= a_\mathcal{I}b_\mathcal{I}\\
	c &\leftarrow \piSOpen{n,2k}(c_\mathcal{I})
\end{align*}

\begin{theorem}
	Let $n \in \N$ and $k \le \frac{n}{2}$. Let $t$ be such that $t \le \frac{n
	- 2k + 1}{2}$ and $t < k$. Then the protocol $\piMulOpen{n,k}$ $t$-securely
	evaluates $\fMulOpen{n,k}$.
\end{theorem}

We require that $k \le \frac{n}{2}$ otherwise after the multiplication the
threshold of the sharing would be too large for it to be possible to
reconstruct the secret even with $n$ parties. The first requirement on $t$
comes from $\piSOpen{n,2k}$, and the second requirement on $t$ is the usual one
to disallow the adversary to reconstruct secrets.

\subsection{Multiplication}

The ideal functionality $\fMul{n,k}$ that represents the multiplication
protocol takes as input two sets of shares for two field elements and gives as
output to each party a share of the product of these two secrets. This
protocol, along with addition, allows for in theory arbitrary arithmetic
circuits and hence arbitrary computations. However, for complicated
computations it is usually more efficient to use other more tailored protocols
and primitives, as we will see is the case for ECDSA\@. The multiplication here
protocol is included mainly for completeness.

Let $a_\mathcal{I}$ and $b_\mathcal{I}$ be two $k$-sharings of respective field
elements $a, b \in \F$. The protocol $\piMul{n,k}(a_\mathcal{I},
b_\mathcal{I})$ is defined by
\begin{align*}
	r_\mathcal{I} &\leftarrow \piRNG{n,k}\\
	c_\mathcal{I} &= a_\mathcal{I}b_\mathcal{I}\\
	c'_\mathcal{I} &= c_\mathcal{I} + r_\mathcal{I}\\
	c' &\leftarrow \piSOpen{n,2k}(c'_\mathcal{I})\\
	c_\mathcal{I} &= c' - r_\mathcal{I}
\end{align*}

\begin{theorem}
	Let $n \in \N$ and $k \le \frac{n}{2}$. Let $t$ be such that $t \le \frac{n
	- 2k + 1}{2}$ and $t < k$. Then the protocol $\piMul{n,k}$ $t$-securely
	evaluates $\fMul{n,k}$.
\end{theorem}

The restrictions on $k$ and $t$ are as for $\piMulOpen{n,k}$. The only value
that we open is $c_\mathcal{I} + r_\mathcal{I}$, which is uniformly random and
so preserves our desired security. Note that $r_\mathcal{I}$ has threshold $k$
whereas $c_\mathcal{I}$ has threshold $2k$. This is not a problem as when we
add two sharings of different thresholds, the resulting sharing has a threshold
equal to the larger of the two. Additionally, we do not need to worry about any
structure this may introduce in the shares, because the sharing is randomised
by $\piSOpen{n,2k}$.

\subsection{Inversion}

The ideal functionality $\fInv{n,k}$ that represents the field inversion
protocol takes as input a sets of shares of a field element and gives as output
to each party shares of the multiplicative inverse of this field element.

Let $a_\mathcal{I}$ be a $k$-sharing of a field element $a \in \F$. The
protocol $\piInv{n,k}(a_\mathcal{I})$ is defined by
\begin{align*}
	r_\mathcal{I} &\leftarrow \piRNG{n,k}\\
	t &\leftarrow \piMulOpen{n,k}(a_\mathcal{I}, r_\mathcal{I})\\
	b_\mathcal{I} &= t^{-1}r_\mathcal{I}
\end{align*}

\begin{remark}
	Because the shares of the product are randomised by $\piMulOpen{n,k}$, it
	is conjectured that it is safe to define $r_\mathcal{I}$ using the less
	expensive biased RNG protocol $\piBRNG{n,k}$.
\end{remark}

\begin{theorem}
	Let $n \in \N$ and $k \le \frac{n}{2}$. Let $t$ be such that $t \le \frac{n
	- 2k + 1}{2}$ and $t < k$. Let $x_\mathcal{I}$ be a $k$-sharing of some $x
	\in \F \setminus \{0\}$. Then the protocol $\piInv{n,k}(x_\mathcal{I})$
	$t$-securely evaluates $\fInv{n,k}(x_\mathcal{I})$.
\end{theorem}

The restrictions on $k$ and $t$ are due to $\piMulOpen{n,k}$. The only value
that is opened is the product of the secret and a random value, and so is
uniformly random. Note however that we need to require that the input shares
are of a non-zero field element, otherwise the open will reveal this fact. Also
note that it is possible that the random number is itself zero, and so if we
were to be careful we would check to see if the value we open is zero, and if
it is, abort this attempt and retry the protocol with a different random
number. However, in our case we ignore this because we will use a field with
approximately $2^{256}$ elements and so the probability that the random number
is zero is negligible.

\section{Signature Algorithm}

The ideal functionality $\fSign{n,k}$ that represents the ECDSA protocol takes
as input a sets of shares of a field element (the private key) and a public
field element (the message digest) and gives as output to each party the values
$r, s$ which constitute a valid ECDSA signature for the given private key and
message.

Let $d_\mathcal{I}$ be a $k$-sharing of a field element $d \in \F$ that
represents an ECDSA private key, and let $z \in \F$ be a message digest. Let
the associated group in which the public keys live, $\G$, have prime order $q$.
The protocol $\piSign{n,k}(d_\mathcal{I}, z)$ is defined by
\begin{align*}
	(k_\mathcal{I}, p_b) &\leftarrow \piRKPG{n,k}\\
	(x, y) &= p_b\\
	r &= x \mod q\\
	k'_\mathcal{I} &\leftarrow \piInv{n,k}(k_\mathcal{I})\\
	t_\mathcal{I} &= rd_\mathcal{I}\\
	t_\mathcal{I} &= t_\mathcal{I} + z_\mathcal{I}\\
	s &\leftarrow \piMulOpen{n,k}(t_\mathcal{I}, k'_\mathcal{I})
\end{align*}

\begin{theorem}
	Let $n \in \N$ and $k \le \frac{n}{2}$. Let $t$ be such that $t \le \frac{n
	- 2k + 1}{2}$ and $t < k$. Let $d_\mathcal{I}$ be a $k$-sharing of a ECDSA
	private key, and let $z \in \F$ be a message digest. Then the protocol
	$\piSign{n,k}(d_\mathcal{I}, z)$ $t$-securely evaluates
	$\fSign{n,k}(d_\mathcal{I}, z)$.
\end{theorem}
The restrictions on $k$ and $t$ are inherited from the constituent protocols.

\begin{remark}
	For the system of interest, we set $k = \frac{n}{3}$. In this case, the
	requirement on $t$ is that $t \le \frac{n}{6} + \frac{1}{2}$. Note however
	that from the discussion of $\piSOpen{n,k}$ that we still maintain
	\textit{safety with abort} for $t \le \frac{n}{3}$.
\end{remark}

\newpage
{%
	% To avoid overfull hboxes in the bibliography
	\emergencystretch=1em

	\printbibliography{}
}

\newpage
\appendix

\section{Measure and Probability Theory}\label{app:prob}

We will use the measure theoretic formulation of probability theory. The key
structure from measure theory is that of a \textit{measure space}, which is
defined as follows.

\begin{definition}
	A measure space is a triple $(X, \Sigma, \mu)$ where $X$ is a set, $\Sigma$
	is a $\sigma$-algebra on $X$ and $\mu$ is a function from $\Sigma$ to
	$\mathbb{R}$ such that the following conditions hold:
	\begin{itemize}
		\item For all $E \in \Sigma$, $\mu(E) \ge 0$
		\item $\mu(\emptyset) = 0$
		\item For any countable collection of disjoint sets $E_1, E_2, \ldots
			\in \Sigma$ we have
			\begin{align*}
				\mu\left(\bigcup_{i=1}^\infty E_i\right) =
				\sum_{i=1}^\infty \mu(E_i)
			\end{align*}
	\end{itemize}
\end{definition}

In this formulation of probability theory, a \textit{probability space} is
simply a special case of a measure space:

\begin{definition}
	A probability space is a measure space $(\Omega, \Sigma, \P)$ such that the
	measure of the entire space is 1; that is,
	\begin{align*}
		\P(\Omega) = 1.
	\end{align*}
\end{definition}

We will often consider the uniform probability measure on a finite set. We
define the standard notation that we will use and formalise the concept as
follows.

\begin{definition}[Uniform probability measure]
	Let $E$ be a finite set. We denote the uniform probability measure on the
	measurable space $(E, \powerset(E))$ by
	\begin{align*}
		U_E: \powerset(E) \to \mathbb{R}.
	\end{align*}
	That is, for each $A \in \powerset(E)$ we have $U_E(A) = \frac{|A|}{|E|}$
	where the notation $|\cdot|$ denotes the number of elements in the given
	set.
\end{definition}

The next key object from probability theory is the \textit{random variable}.
This is formalised as a measurable function.

\begin{definition}
	Let $(\Omega, \Sigma, \P)$ be a probability space and $(E, \mathcal{E})$ be
	a measurable space. A random variable $X$ is a measurable function $X:
	\Omega \to E$.
\end{definition}

Note that the random variable $X$ induces a probability measure $\mu$ on the
measurable space $(E, \mathcal{E})$ which we can define by
\begin{align*}
	\mu: \mathcal{E} &\to \mathbb{R}\\
	A &\mapsto \P(\set{\omega \in \Omega}{X(\omega) \in A}),
\end{align*}
and so we also have the induced probability space $(E, \mathcal{E}, \mu)$. Also
note that for a measurable space $(F, \mathcal{F})$ and measurable function $f:
E \to F$, we can define a new random variable $Y = f \circ X$. Thus when we
write $f(X)$ we understand this to be the random variable $Y$. It is also
useful to define the following common short hand notations:

\begin{definition}
	Let $(\Omega, \Sigma, \P)$ be a probability space, $(E, \mathcal{E})$ be a
	measurable space and $X: \Omega \to E$ be a random variable. Let the
	induced probability measure on $(E, \mathcal{E})$ be $\mu$. Define the
	following short hand notations:
	\begin{itemize}
		\item $\P(X = x) = \mu(X = x) = \P(\set{\omega \in
			\Omega}{X(\omega) = x}) \quad \forall x \in E$
		\item $\P(X \in A) = \mu(X \in A) = \mu(A) \quad \forall A \in
			\mathcal{E}$
	\end{itemize}
\end{definition}

While the latter is actually not shorter than the technically correct version,
it is conceptually clearer which is why we use it.

\section{Proof of Theorem~\ref{thm:unif}}\label{app:unif}

The following is the proof of Theorem~\ref{thm:unif}

\begin{proof}
	We need only show that $\mu(X = e) = \frac{1}{n}$ for all $e \in E$. To do
	this, define
	\begin{align*}
		A_{i,e} = \set{\omega \in \Omega}{f_i(\omega) = e}.
	\end{align*}
	Since each $f_i$ is measurable, it follows that $A_{i,e} \in \Sigma$ for
	all $i \in \seq{n}$ and $e \in E$. Notice that these sets partition
	$\Omega$; this is summarised in the following claim.
	\begin{claim}
		The sets $A_{i,e}$ for $i \in \seq{n}$ partition $\Omega$.
	\end{claim}
	\begin{proof}
		To show this, we need to show that each $A_{i,e}$ is disjoint, and that
		they cover $\Omega$. The first property holds due to the assumption
		that $f_i(\omega) \ne f_j(\omega)$ for all $i \ne j$, $i, j \in
		\seq{n}$ and for all $\omega \in \Omega$. This assumption also shows
		that for all $\omega \in \Omega$ there must exist some $i \in \seq{n}$
		for which $f_i(\omega) = e$, as $E$ only has $n$ elements. But this is
		precisely the statement that $\omega \in A_{i,e}$. This completes the
		proof.
	\end{proof}
	Also, these sets define a partition of our set of interest:
	\begin{align*}
		\bigcup_{i=1}^n A_{i,e} \times \{f_i\}
		= \set{(\omega, f) \in \Omega \times F}{X(\omega, f) = e}.
	\end{align*}
	Thus for any $e \in E$ we have
	\begin{align}
		\label{eq:unif_thm}\mu(X = e) =
		\sum_{i=1}^n \mu\left(A_{i,e} \times \{f_i\}\right).
	\end{align}
	By the definition of the product measure we know that
	\begin{align*}
		\mu(A_{i,e} \times \{f_i\}) = \P(A_{i,e})U_F(\{f_i\}),
	\end{align*}
	and so using this and the fact that $A_{1,e}, \ldots, A_{n,e}$ partitions
	$\Omega$ we can reduce~\eqref{eq:unif_thm} to
	\begin{align*}
		\mu(X = e) &= \frac{1}{n}\sum_{i=1}^n \P(A_{i,e})\\
		&= \frac{1}{n},
	\end{align*}
	which is the desired result.
\end{proof}

\section{Shamir Secret Sharing Theorem Proofs}\label{app:shamir}

First we present the proof of Theorem~\ref{thm:shareDist}.
\begin{proof}
	We consider three separate cases.
	\begin{itemize}
		\item $t \ge k$: In this case, we can view
			equation~\eqref{eq:shareDist} as a set of $t$ linear equations with
			$k$ unknowns, and hence has a unique solution for $(x, c_1, \ldots,
			c_{k-1})$ (and in the cases where $t > k$ and the set of equations
			is overconstrained, we know that it will not have no solutions
			because we have defined the sharing to be correct). Let this
			solution be $z = (y, z_1, \ldots, z_{k-1})$. Then from independence
			we have
			\begin{align*}
				\P_t(X = z) &=
				\P_t\left(
					x = y \cap \bigcap_{i = 1}^{k-1} c_i = z_i
				\right)\\
				&= \P_t(x = y)
				\P_t((c_1, \ldots, c_{k-1}) = (z_1, \ldots, z_{k-1}))\\
				&= \mu(\{y\})|\F|^{-k+1},
			\end{align*}
			which is the desired result for this case.

		\item $t = k - 1$: Suppose that we fix $x$ in
			equation~\eqref{eq:shareDist}. We will again have a system of (in
			this case $k-1$) linear equations with as many unknowns, and so we
			must have a unique solution $d_y$ for $(c_1, \ldots, c_{k-1})$ for
			a given set of shares $z$. We may therefore write
			\begin{align*}
				\P_t(X = z) &=
				\sum_{y \in \F}
				\P_t(x = y \cap (c_1, \ldots, c_{k-1}) = d_y)\\
				&= \sum_{y \in \F} \mu(\{y\})|\F|^{-k+1}\\
				&= |\F|^{-k+1}.
			\end{align*}
			Since $t = k - 1$ this is the desired result.

		\item $t < k - 1$: Here we take a similar approach to the previous
			case. If we fix $x = y$ but also $(c_{t+1}, \ldots, c_{k-1}) = c$
			in equation~\eqref{eq:shareDist}, then we have $t$ linear equations
			in the $t$ unknowns $c_1, \ldots, c_t$, which has a unique solution
			for a given set of shares $z$ that we will denote $d_{y, c}$. We
			may then write the probability $P(X = z)$ as
			\begin{align*}
				\sum_{y \in \F}\sum_{c \in \F^{k - 1 - t}}
				\P_t(
					x = y \cap
					(c_1, \ldots, c_t) = d_{y, c} \cap
					(c_{t+1}, \ldots, c_{k-1}) = c
				),
			\end{align*}
			which we can again use independence and our known distributions to
			simplify to
			\begin{align*}
				\sum_{y \in \F}\sum_{c \in \F^{k - 1 - t}}
				\mu(\{y\})|\F|^{-t}|\F|^{-k+1+t}
				=
				|\F|^{-t}.
			\end{align*}
			This final case completes the proof.
	\end{itemize}
\end{proof}

Now we present a proof of Theorem~\ref{thm:shareCoeffLO}.

\begin{proof}
	Consider the \textit{Vandermonde} matrix for the elements ${(x_i)}_{i \in
	\mathcal{R}}$, which is defined as
	\begin{align*}
		V =
		\begin{pmatrix}
			1 & i_1 & i_1^2 & \cdots & i_1^t\\
			1 & i_2 & i_2^2 & \cdots & i_2^t\\
			1 & i_3 & i_3^2 & \cdots & i_3^t\\
			\vdots & \vdots & \vdots & \ddots & \vdots\\
			1 & i_k & i_k^2 & \cdots & i_k^t\\
		\end{pmatrix},
	\end{align*}
	where $\mathcal{R} = \{i_1, \ldots, i_k\}$. If we define $s$ to be the
	$k$-vector of shares ${(x_i)}_{i \in \mathcal{R}}$, and $c$ to be the $k$
	vector of coefficients, then by the way secret sharing is defined we have
	the relationship
	\begin{align*}
		Vc = s.
	\end{align*}
	It is well known that the determinant of the Vandermonde matrix $V$ above
	is
	\begin{align*}
		\det(V) = \prod_{1 \le p < q \le k} (i_p - i_q),
	\end{align*}
	and is hence non-zero, if $i_p \ne i_q$ for all $p,q \in \mathcal{R}$, $p
	\ne q$, which in our case is true. This means that $V$ is invertible, and
	so we may write
	\begin{align*}
		c = V^{-1}s.
	\end{align*}
	We can thus complete the proof by defining $\lambda_i^{(j)}$ as the $(j,
	i)$ element of $V^{-1}$.
\end{proof}

\section{Proof of security of $\piSOpen{n,k}$}\label{app:so}

Here we provide full mathematical details for the proof of
Theorem~\ref{thm:so}.

\begin{proof}
	\newcommand{\proto}{\piSOpen{n,k}}

	Let an adversary $\mathcal{A}$ be given. Let input $x$ and auxiliary input
	$z$ be given. We will outline briefly our proof strategy, as we will also
	use it in other proofs. By definition, our end goal is to construct
	$\mathcal{S}$ such that
	\begin{align*}
		\execD \eqd \idealD
	\end{align*}
	We could try to compare these random variables directly, but in general
	this is a little cumbersome and complicated, as many elements of these
	random variables will have dependencies on each other. In addition, the
	view for $\mathcal{S}$ will have to be constructed by reverse-engineering
	based on the output of the trusted party, which makes the equality of the
	distributions less clear. To try to make things easier, we will identify
	the key source of randomness that determines these views, and a
	deterministic function that maps this to the view. For example, if a view
	contains all of the shares of some $k$-sharing, instead of considering the
	shares directly, we could use only the $k$ coefficients that determine the
	sharing, and then we can map these deterministically to the shares.
	Specifically, we seek random variables $X$ and $Y$ (the simplified
	randomness) and a deterministic function $h$ such that $h(X) \eqd \execD$
	and $h(Y) \eqd \idealD$. Then, we will show that $X \eqd Y$, at which point
	we may apply Theorem~\ref{thm:eqd_det} to arrive at our desired result.

	In our current case, we begin by characterising $\execD$. The first
	messages that $\mathcal{A}$ receives are the $t$ output messages
	${(z_i)}_{i \in \mathcal{I}_c}$ from $\piRZG{n, k}$. The next part of the
	view comes from the $n - t$ messages ${(y_i)}_{i \in \mathcal{I}_h}$ that
	the uncorrupted parties send in the broadcast round. Next, $\mathcal{A}$
	sends its $t$ messages ${(m_i)}_{i \in \mathcal{I}_c}$ in the broadcast
	round. The final part of $\execD$ is the outputs of the parties; this is
	$t$ $\bot$s that are output by the corrupted parties and $n-t$ elements of
	$\F$ (that should all be the same and equal to the secret corresponding to
	the sharing $x$). We can thus write
	\begin{align*}
		\execD
		=
		\left(
			{(z_i)}_{i \in \mathcal{I}_c},
			{(y_i)}_{i \in \mathcal{I}_h},
			{(m_i)}_{i \in \mathcal{I}_c},
			s'
		\right),
	\end{align*}
	where for simplicity we ignore the outputs of the corrupted parties and
	only include one field element $s'$ to represent the output of the honest
	parties.

	The ``key randomness'' for $\execD$ as discussed at the beginning of this
	proof is captured by the random variable
	\begin{align*}
		X
		=
		\left(
			{(c_i)}_{i \in \seqZ{k-1}},
			{(m_i)}_{i \in \mathcal{I}_c}
		\right),
	\end{align*}
	where ${(c_i)}_{i \in \seqZ{k-1}}$ are independently and uniformly randomly
	distributed elements of $\F$, representing the coefficients for the sharing
	${(z_i)}_{i \in \mathcal{I}}$ of zero, and ${(m_i)}_{i \in \mathcal{I}_c}$
	are the messages that $\mathcal{A}$ would send in the broadcast round. We
	describe the latter more specifically. The messages that $\mathcal{A}$ will
	send in any given round depends on its random tape and all of the messages
	it has received up to that point. The random tape is picked uniformly
	randomly, so we need only consider the messages it has received up until
	the broadcast round; these are its shares of zero and the randomised shares
	sent by the honest parties. We will assume that the shares of zero come
	from a uniformly random sharing of zero, and that the shares received from
	the honest parties come from a uniformly random sharing of $x$, the secret
	corresponding to the input shares. If we let the coefficients that define
	the input sharing be $c_x \in \F^k$, then we can define our deterministic
	function $h$ by the mapping
	\begin{align*}
		(c, m)
		\mapsto
		\left(
			\spl(\mathcal{I}_c, c),
			\spl(\mathcal{I}_h, c + c_x),
			m,
			s
		\right),
	\end{align*}
	where we interpret $c + c_x$ to be element-wise addition.

	We want to show that $h(X) \eqd \execD$. As previously mentioned, the third
	element (the messages sent by $\mathcal{A}$) depend only on the previous
	two elements, as these are the messages received up until the broadcast
	round. This means that if we ensure that these first two elements have the
	correct distribution, the third will too. Now, notice that the final
	element of $\execD$ will be fixed and equal to $s$. This is because by
	definition and the assumptions on $t$ the honest parties will hold at least
	$\frac{n + k - 1}{2}$ correct shares and so can always correctly
	reconstruct $s$. Since this final element is fixed for the given inputs,
	and in $h$ is defined to be that same fixed value $s$, the final element
	will always be correct. Finally, we notice that in fact the first element
	is a deterministic function of the second element. This is because since
	$|\mathcal{I}_h| \ge k$ we know that ${(y_i)}_{i \in \mathcal{I}_h}$
	completely determines the rest of the shares ${(y_i)}_{i \in
	\mathcal{I}_c}$ and the random shares of zero satisfy $z_i = y_i - x_i$ for
	all $i \in \mathcal{I}_c$, where $x = {(x_i)}_{i \in \mathcal{I}}$. Thus we
	consider only the second element. In $\execD$, this element is distributed
	as a part of a uniformly random sharing of $s$. This follows from the fact
	that the output shares of $\piRZG{n,k}$ will be a uniformly random sharing
	of zero, and Corollary~\ref{cor:unif}. In $h(X)$, this element will also
	have this distribution for exactly the same reason; $c$ is distributed in
	the same way as the coefficients of a uniformly random sharing of zero. We
	therefore conclude that $h(X) \eqd \execD$.

	With this result established, we now seek to construct $\mathcal{S}$ and a
	random variable $Y$ such that $h(Y) \eqd \idealD$ and $X \eqd Y$. Define
	the simulator adversary $\mathcal{S}$ as follows.

	\begin{enumerate}
		\item\label{enum:openGen} Let the output from the trusted party be
			$s$ (by definition it will always be the secret corresponding to
			the sharing $x$). $\mathcal{S}$ begins by constructing a random
			sharing ${(y_i)}_{i \in \mathcal{I}}$ of $s$.

		\item $\mathcal{S}$ now defines
			\begin{align*}
				z_i = y_i - x_i \quad \forall i \in \mathcal{I}_c.
			\end{align*}

		\item\label{enum:openMsg} $\mathcal{S}$ now runs $\mathcal{A}$ after
			giving it ${(z_i)}_{i \in \mathcal{I}_h}$ and ${(y_i)}_{i \in
			\mathcal{I}_h}$ to get the messages ${(m_i)}_{i \in \mathcal{I}_c}$
			that $\mathcal{A}$ sends in the broadcast round.

		\item $\mathcal{S}$ then outputs
			\begin{align*}
				{(z_i)}_{i \in \mathcal{I}_c},
				{(m_i)}_{i \in \mathcal{I}_c},
				{(y_i)}_{i \in \mathcal{I}_h}
			\end{align*}
			as the generated view.
	\end{enumerate}

	We now aim to construct $Y$. Let ${(y_i)}_{i \in \mathcal{I}}$ and
	${(m_i)}_{i \in \mathcal{I}_c}$ be distributed as in the description of
	$\mathcal{S}$. Let the coefficients corresponding to this sharing be $c_y =
	{(c_i^{(y)})}_{i \in \seqZ{k-1}}$ and let $c_x = {(c_i^{(x)})}_{i \in
	\seqZ{k-1}}$, and define the coefficients $c = {(c_i)}_{i \in \seqZ{k-1}}$
	by
	\begin{align}\label{eq:sopenCoeff}
		c_i = c_i^{(y)} - c_i^{(x)} \quad \forall i \in \seqZ{k-1}.
	\end{align}
	We then define
	\begin{align*}
		Y
		=
		\left(
			{(c_i)}_{i \in \seqZ{k-1}},
			{(m_i)}_{i \in \mathcal{I}_c}
		\right).
	\end{align*}
	With $Y$ defined, our next step is to show that $h(Y) \eqd \idealD$. The
	same argument as before shows that to do this we need only consider the
	second element of these distributions. For this element, the honest shares
	${(y_i)}_{i \in \mathcal{I}_c}$, these clearly have the same distribution
	since $c + c_x = c_y$ by definition and $\spl(\mathcal{I}_h, c_y)$ has the
	same distribution as the corresponding element in $\idealD$ also by
	definition. Thus $h(Y) \eqd \idealD$.

	The final step for our proof is to show that $X \eqd Y$. The first element
	in $X$ is distributed as $k-1$ independently and uniformly randomly
	distributed elements of $\F$ (the coefficient $c_0$ is $0 \in \F$ as it is
	a sharing of zero). In $Y$ the distribution is the same, which follows from
	Eq.~\eqref{eq:sopenCoeff}, Corollary~\ref{cor:unif} and the fact that each
	$c_i^{(y)}$ is independently and uniformly distributed for each $i \in
	\seq{k-1}$, and also the fact that $c_0^{(y)} = c_0^{(x)}$. Finally, the
	second elements have the same distribution due to their construction and
	the fact that they otherwise only depend on the first element.
\end{proof}

\section{Proof of Security of $\piRNG{n,k}$}\label{app:rng}

Here we provide full mathematical details for the proof of
Theorem~\ref{thm:rng}.

\begin{proof}
	\newcommand{\proto}{\piRNG{n,k}}

	Let adversary $\mathcal{A}$, input $x$ and auxiliary input $z$ be given.

	We will characterise $\execD$. Define $m$ to be the messages that were sent
	and received during step~\ref{enum:rngCoeff} (except for the output shares
	of each invocation of $\piBRNG{n,k}$, these will be labelled separately)
	and also define $m_o$ to be the messages sent by $\mathcal{A}$ for each of
	the invocations of $\piDOpen{n,k,j}$ in step~\ref{enum:rngDO}. It follows
	that $\execD$ is equal to
	\begin{align*}
		m,
		{\left(
			{\left(r_i^{(j)}\right)}_{i \in \mathcal{I}_c}
		\right)}_{j \in \seq{k}},
		m_o,
		{\left(
			r_i, {(r_{j,i})}_{j \in \mathcal{I}}
		\right)}_{i \in \mathcal{I}_c},
		{(y_i)}_{i \in \seq{n}},
	\end{align*}
	where $y_i$ is defined as in $\piRNG{n,k}$ if $i \in \mathcal{I}_h$, and
	$\bot$ otherwise. Note that this is independent of the input $x$ to
	$\piRNG{n,k}$ as any such input is ignored by the protocol (and also by
	$\fRNG{n,k}$). Denote the set in which this lives by $F$; i.e. $\execD \in
	F$.

	We now seek to define a random variable $X$ taking values in some set $E$
	and a function $h: E \to F$ in order to eventually apply
	Theorem~\ref{thm:eqd_det}. To do this, first define the random variable $X$
	as
	\begin{align*}
		X =
		\left(
			m,
			{\left(
				{\left(r_i^{(j)}\right)}_{i \in \mathcal{I}}
			\right)}_{j \in \seqZ{k-1}},
			m_o
		\right),
	\end{align*}
	which is produced by $\mathcal{A}$ interacting with the trusted party for
	$\piBRNG{n,k}$, and where the random tape for $\mathcal{A}$ and the trusted
	party is uniformly randomly chosen. Define the set $E$ implicitly as the
	set that $X$ takes values in. Now we can define our function $h$ as
	\begin{align*}
		\begin{aligned}
			h: E &\to F\\
			\left(
				m,
				r,
				m_o
			\right)
			&\mapsto
			\left(
				m,
				\nu(r),
				m_o,
				{\left(
					\spl(\{i\}, \mu(r)),
					{\left(\varphi(r, j, i)\right)}_{j \in \mathcal{I}}
				\right)}_{i \in \mathcal{I}_c},
				{\left(\lambda(i, \mu(r))\right)}_{i \in \mathcal{I}}
			\right),
		\end{aligned}
	\end{align*}
	where we have the following definitions:
	\begin{itemize}
		\item $\nu$ simply projects $r$, which includes shares for all parties,
			to just those shares that the corrupted parties receive; i.e.
			\begin{align*}
				\nu
				\left(
					{\left(
						{\left(
							r_i^{(j)}
						\right)}_{i \in \mathcal{I}}
					\right)}_{j \in \seqZ{k-1}}
				\right)
				=
				{\left(
					{\left(r_i^{(j)}\right)}_{i \in \mathcal{I}_c}
				\right)}_{j \in \seqZ{k-1}}.
			\end{align*}

		\item $\mu$ maps the set of shares
			\begin{align*}
				{\left(
					{\left(r_i^{(j)}\right)}_{i \in \mathcal{I}}
				\right)}_{j \in \seqZ{k-1}}
			\end{align*}
			to their corresponding (uniquely defined) secrets for each $j \in
			\seqZ{k-1}$; the output is ${(c_i)}_{i \in \seqZ{k-1}}$. In the
			protocol these are the coefficients corresponding to the final
			random sharing.

		\item $\varphi$ converts the shares of the coefficients into shares of
			the shares, as in step~\ref{enum:rngrij} of the protocol. It is
			defined as
			\begin{align*}
				\varphi: \F^{nk} \times \F^2
				&\to
				\F^n\\
				{\left(
					{\left(
						r_i^{(j)}
					\right)}_{i \in \mathcal{I}}
				\right)}_{j \in \seqZ{k-1}},
				(i, j)
				&\mapsto
				\sum_{l=0}^{k-1} r_i^{(l)} j^l
			\end{align*}

		\item $\lambda$ is defined by
			\begin{align*}
				\lambda: \F \times \F^k
				&\to
				\F \cup \{\bot\}\\
				\left(
					i,
					{(c_j)}_{j \in \seqZ{k-1}}
				\right)
				&\mapsto
				\begin{cases}
					\bot & i \in \mathcal{I}_c\\
					\sum_{j=0}^{k-1} c_j i^j & i \notin \mathcal{I}_c
				\end{cases},
			\end{align*}
	\end{itemize}

	Now, we claim that $\execD \eqd h(X)$. This is immediate for the first
	three elements, as they are generated in $X$ exactly as they are in
	$\piRNG{n,k}$ and using the same distributions. As for the second last
	element, which represents the messages during the invocations of
	$\piDOpen{n,k,j}$, this is correct due to how it is defined from $r$ in
	$h(X)$ and the correctness property that the output of $\piDOpen{n,k}$
	enjoys from its proof of security. Similarly, the fact that the last
	element has the correct distribution also follows from its definition in
	$\execD$ and in $h(X)$.

	Now that we have characterised $\execD$, we turn our attention to the
	simulated view. The idea behind the simulator is as follows. The simulator
	can run the protocol as usual up until step~\ref{enum:rngDO}. At this
	point, $\mathcal{A}$ does not know what the secrets for the random sharings
	are, and the honest parties shares will be uniformly random conditioned on
	the fact that they are consistent with $\mathcal{A}$'s shares. This allows
	the simulator to simply extend the corrupted parties' shares (which are
	known by the simulator since it holds enough shares to reconstruct them) to
	valid random sharings of the target output shares, and this will have the
	correct distribution.

	The simulator is defined as follows.
	\begin{enumerate}
		\item The trusted party is invoked first. Let the output shares of the
			trusted party be ${(y_i)}_{i \in \mathcal{I}}$.

		\item $\mathcal{S}$ runs $\mathcal{A}$ and acts on behalf of the honest
			parties for Step~\ref{enum:rngCoeff}. In doing so, $\mathcal{S}$
			learns the biased shares ${\left(r_i^{(j)}\right)}_{j \in
			\seqZ{k-1}}$ for each $i \in \mathcal{I}_c$, i.e.\ for each
			corrupted party. Label the messages sent and received by
			$\mathcal{A}$ during the invocations of $\piBRNG{n,k}$ as $m$, and
			the messages that $\mathcal{A}$ would send to the invocations of
			$\piDOpen{n,k,j}$ as $m_o$.

		\item $\mathcal{S}$ constructs the honest parties' shares
			${\left({\left(r_i^{(j)}\right)}_{i \in \mathcal{I}_h}\right)}_{j
			\in \seqZ{k-1}}$ as follows. First, $\mathcal{S}$ picks random
			coefficients ${(c_i)}_{i \in \seqZ{k-1}}$ that are consistent with
			the target shares ${(y_i)}_{i \in \mathcal{I}_c}$, i.e.\
			$\mathcal{S}$ picks uniformly randomly from the set
			\begin{align*}
				\set{%
					{(c_i)}_{i \in \seqZ{k-1}} \in \F^k
				}{%
					y_i = \sum_{j=0}^{k-1} c_j i^j
					\;
					\forall i \in \mathcal{I}_c
				}.
			\end{align*}
			Next, the honest parties' shares are chosen at random under the
			condition that ${\left(r_i^{(j)}\right)}_{i \in \mathcal{I}}$ is a
			consistent sharing of $c_j$ for each $j \in \seqZ{k-1}$. Precisely,
			the shares are chosen uniformly randomly from the set
			\begin{align*}
				\set{%
					{\left(
						{\left(
							r_i^{(j)}
						\right)}_{i \in \mathcal{I}_h}\right)
					}_{j \in \seqZ{k-1}}
					\in
					\F^{k(n-t)}
				}{%
					\begin{matrix}
						\forall j \in \seqZ{k-1}\\
						\exists {(c_l')}_{l \in \seq{k-1}}\\
						\forall i \in \mathcal{I}
					\end{matrix}
					:
					r_i^{(j)} = c_j + \sum_{l=1}^{k-1} c_l' i^l
				}.
			\end{align*}

		\item Let the random variable $Y$ be defined as
			\begin{align*}
				Y =
				\left(
					m,
					{\left(
						{\left(r_i^{(j)}\right)}_{i \in \mathcal{I}}
					\right)}_{j \in \seqZ{k-1}},
					m_o
				\right).
			\end{align*}
			$\mathcal{S}$ finishes by constructing its output by computing
			$h(Y)$ and discarding the last element, which corresponds to the
			output of the protocol; recall that this is not part of the
			simulator's output as instead this is defined by the output of the
			trusted party.
	\end{enumerate}

	Now we need to consider $\idealD$. We want to show that $h(Y) \eqd
	\idealD$. This is clearly true for all but the last element because of how
	$\idealD$ was constructed. But the last element will also have the correct
	distribution with regard to the others due to its construction, as $Y$ was
	reverse engineered to be consistent with the given output of the trusted
	party.

	Finally, we want to show that $X \eqd Y$, as then by
	Theorem~\ref{thm:eqd_det} we will have $h(X) \eqd h(Y)$, which gives the
	last equality needed to show that
	\begin{align*}
		\execD \eqd \idealD
	\end{align*}
	which will complete the proof. By construction, this is clearly true for
	all but the honest party shares
	\begin{align*}
		{\left(
			{\left(r_i^{(j)}\right)}_{i \in \mathcal{I}_h}
		\right)}_{j \in \seqZ{k-1}},
	\end{align*}
	so we need only focus our attention on these. In the case of $X$, we know
	from the secrecy property of the output shares of $\piBRNG{n,k}$ that
	\begin{align*}
		r_h^{(j)}
		=
		{\left(r_i^{(j)}\right)}_{i \in \mathcal{I}_h}
	\end{align*}
	is independently and uniformly distributed in
	$\shareRestrict{r_c^{(j)}}{\mathcal{I}}{\mathcal{I}_c}$ for each $j \in
	\seqZ{k-1}$, where
	\begin{align*}
		r_c^{(j)} = {\left(r_i^{(j)}\right)}_{i \in \mathcal{I}_c}.
	\end{align*}
	The result for $Y$ follows from the proceeding claims, the proofs of which
	will conclude the proof of security.

	\begin{claim}
		The coefficients $c = c_0, \ldots, c_{k-1}$ as defined in the
		simulation are independently and uniformly randomly distributed.
	\end{claim}

	\begin{proof}
		Since the coefficients of a $k-1$ degree polynomial are completely
		determined by $k$ points on that polynomial, we can see that for a
		given set of shares $y = {(y_i)}_{i \in \mathcal{I}_c}$ (which is $t$
		points), we could choose any set of $k-t$ points to determine the
		coefficients. Thus for a given $y$ there are $|\F|^{k-t}$ possible
		choices for $c$, each one being picked with equal probability. Now,
		consider the probability of obtaining some given $c$, using the fact
		that $y$ is chosen uniformly randomly. For this given $c$, there is a
		unique $y$ that is consistent with $c$, and this $y$ is chosen with
		probability $|\F|^{-t}$. But we not only require that the uniquely
		determined $y$ was chosen, but also the unique $c$ for the
		possibilities determined by $y$, for which we have just seen there
		exists $|\F|^{k-t}$ choices, each chosen with equal probability.
		Putting these two results together, we find that the probability for a
		given $c$ is $|\F|^{-k}$. Since $c \in \F^k$, the result follows.
	\end{proof}

	\begin{claim}
		For each $j \in \seqZ{k-1}$, the shares $r_h^{(j)}$ as defined in the
		simulation are independently and uniformly randomly distributed in
		$\shareRestrict{r_c^{(j)}}{\mathcal{I}}{\mathcal{I}_c}$.
	\end{claim}

	\begin{proof}
		The independence of each $r_h^{(j)}$ follows easily from the
		independence of each $c_j$. Next, notice that for a given $c_j$,
		$r_h^{(j)} \in A_{c_j}$ where for each $c \in \F$ we have
		\begin{align*}
			A_c
			=
			\set{%
				r_h^{(j)} \in \F^{n-t}
			}{%
				\exists {(c_i')}_{i \in \seq{k-1}}:
				\forall i \in \mathcal{I}_c,
				r_i = c + \sum_{l=1}^{k-1} c_l' i^l
			}.
		\end{align*}
		Since a polynomial with degree $k$ is completely determined by $k$
		points, so too are its coefficients, and hence each distinct element
		(which is a set of $n-t > k$ shares) in $A_c$ corresponds also to a
		distinct set of coefficients $c, c_1', \ldots, c_{k-1}'$. From this it
		follows that for $c \ne d$ we have that $A_c$ and $A_d$ are disjoint
		sets. Now, since $r_h^{(j)}$ is chosen uniformly randomly from
		$A_{c_j}$, and $c_j$ is uniformly random, it follows that $r_h^{(j)}$
		is chosen uniformly randomly from the set
		\begin{align*}
			\bigcup_{c_j \in \F} A_{c_j}.
		\end{align*}
		But it is easy to see that this set is equal to
		$\shareRestrict{r_c^{(j)}}{\mathcal{I}}{\mathcal{I}_c}$.
	\end{proof}

	This completes the proof.
\end{proof}

\end{document}
